%% ============================================================================
%% Chapter 6: Building Production REST APIs with FastAPI
%% Mastering APIs: From Zero to Production Guru
%% ============================================================================
\chapter{Building Production REST APIs with FastAPI}
\label{ch:fastapi}

%% ----------------------------------------------------------------------------
%% Executive Summary
%% ----------------------------------------------------------------------------
Chapters 1--5 taught you to speak HTTP fluently and to consume other people's
APIs defensively. This chapter turns the table: you are now the
\emph{provider}. Building an API that strangers can rely on is a different
discipline from writing one that works on your laptop --- the contract must be
explicit, the runtime must survive hostile input and partial failure, the
operational surface (health, logs, metrics, shutdown) must exist \emph{before}
the first deploy, and every one of those properties must be testable without
a network. \textbf{FastAPI}~\cite{fastapidocs} is the framework this book uses
for that job: a thin, ruthlessly pragmatic layer over Starlette (routing,
middleware, the ASGI runtime) and Pydantic~v2 (validation and serialization)
that turns Python type annotations into a machine-readable OpenAPI~3.1
contract~\cite{openapi31}.

We begin where production systems actually begin --- with the contract --- and
confront the design-first versus code-first dilemma head-on, settling on a
disciplined answer: the OpenAPI document is the backbone either way; the
question is only who writes it. We then descend into the runtime: the ASGI
callable protocol (\texttt{scope}/\texttt{receive}/\texttt{send}), what uvicorn
and hypercorn actually do, where Starlette ends and FastAPI begins, and what
the event loop means for your choice between \texttt{def} and
\texttt{async~def} endpoints. The centerpiece is a complete, production-grade
\emph{Northwind Robotics Orders API}: typed Pydantic~v2 models with domain
invariants, CRUD routes, a repository boundary with in-memory and SQLite
implementations behind dependency injection, pure-ASGI middleware carrying
correlation IDs into structured JSON logs, liveness/readiness probes, custom
exception handlers that never leak stack traces, an NDJSON streaming endpoint,
and a customized OpenAPI document --- all verified by an in-process
\texttt{TestClient} suite that runs offline in under a second.

By the end of this chapter you will be able to explain --- at the level of
interview-grade precision --- exactly what happens between a TCP byte stream
arriving at uvicorn and a validated, serialized, logged response leaving it,
and you will be able to build that pipeline yourself.

\begin{keyconceptbox}
FastAPI is not a web framework so much as a \textbf{compiler from type
annotations to contracts}. Declare the shape of your data with Pydantic
models, declare how to obtain your collaborators with \texttt{Depends}, and
declare your endpoint signatures; FastAPI derives validation, serialization,
dependency resolution, and the OpenAPI schema from those declarations. Your
job shifts from writing plumbing to writing \emph{declarations you can
defend}.
\end{keyconceptbox}

%% ----------------------------------------------------------------------------
\section{Learning Objectives \& Prerequisites}
\label{sec:ch6-objectives}

\subsection*{Learning Objectives}
After completing this chapter, its lab, and its exercises, you will be able to:

\begin{enumerate}[label=\textbf{LO-\arabic*.}, leftmargin=2.6cm]
  \item Argue the design-first versus code-first trade-off quantitatively, and
        construct an OpenAPI~3.1 document --- \texttt{servers},
        \texttt{paths}, operations, \texttt{components.schemas}, and
        \texttt{securitySchemes} --- that serves as a source of truth for
        implementation, documentation, and client
        generation~\cite{openapi31}.
  \item Explain the ASGI callable protocol (\texttt{scope},
        \texttt{receive}, \texttt{send}), the division of labor between
        uvicorn/hypercorn, Starlette, and FastAPI, and the event-loop
        implications of every line of handler code you write.
  \item Trace the complete FastAPI request lifecycle --- routing resolution,
        middleware stack, dependency resolution, Pydantic validation, handler
        invocation, response serialization --- and predict the ordering
        guarantees of each stage.
  \item Design dependency-injection graphs with \texttt{Depends}: sub-
        dependencies, per-request caching, \texttt{yield}-based cleanup, and
        \texttt{dependency\_overrides} for tests.
  \item Choose correctly between \texttt{def} and \texttt{async~def}
        endpoints using threadpool semantics, and between
        \texttt{BackgroundTasks} and a real task queue using durability and
        latency criteria.
  \item Apply the Pydantic~v2 validation pipeline ---
        \texttt{field\_validator}, \texttt{model\_validator}, serialization
        aliases, and the Rust-backed core --- to enforce domain invariants,
        and use \texttt{response\_model} as a security boundary that prevents
        internal fields from leaking~\cite{pydanticdocs}.
  \item Operate the service: 12-factor configuration with
        \texttt{pydantic-settings}~\cite{twelvefactor}, structured JSON
        logging with correlation IDs, liveness versus readiness probes, and
        graceful shutdown via the lifespan protocol.
  \item Build, run, and test the complete Orders API offline with
        \texttt{TestClient}/httpx~\cite{httpxdocs}, including 422 validation
        errors, domain errors, health endpoints, and OpenAPI schema
        assertions.
\end{enumerate}

\subsection*{Prerequisites}
This chapter assumes you have completed:

\begin{itemize}
  \item \textbf{Chapter 0} --- environment setup: Python 3.11+, a virtual
        environment, and the ability to run \texttt{pip} and \texttt{pytest}
        from PowerShell.
  \item \textbf{Chapter 1} --- the HTTP wire protocol: methods, status codes,
        header fields, and message bodies. Every route in this chapter is an
        exercise of that vocabulary.
  \item \textbf{Chapter 2} --- consuming APIs from Python: \texttt{httpx}
        basics, which the test suite reuses in its ASGI-transport form.
  \item \textbf{Chapter 3} --- the REST architectural style: resources,
        representations, and the uniform interface that the Orders API
        implements.
  \item \textbf{Chapter 4} --- serialization and data contracts: JSON, NDJSON,
        and the notion of a machine-checkable contract, which OpenAPI
        operationalizes here.
\end{itemize}

No network access is required anywhere in this chapter: every test runs
against the in-process ASGI application, and every dataset is synthetic,
embedded, and deterministic, set in the fictional \emph{Northwind Robotics}
parts and logistics universe.

\begin{warningbox}
\textbf{Version discipline.} The code in this chapter targets
FastAPI~$\geq0.110$, Starlette's modern lifespan API, and Pydantic~v2
($\geq2.7$). Pydantic v1 idioms --- \texttt{@validator}, \texttt{.dict()},
\texttt{class Config} --- appear constantly in old blog posts and in large
language model output; they are \emph{wrong} here and several raise hard
errors under v2. Pin the versions from Section~\ref{sec:ch6-deps} before
writing a line of code.
\end{warningbox}

%% ----------------------------------------------------------------------------
\section{Deep Technical Mechanics \& Architectural Concepts}
\label{sec:ch6-mechanics}

\subsection{Design-First vs.\ Code-First: Who Writes the Contract?}
\label{sec:ch6-design-first}

Every HTTP API has a contract --- a description of its resources, operations,
representations, and failure modes --- whether or not anyone writes it down.
The only choice is whether that contract is an \emph{artifact you author} or
an \emph{archaeological finding} your consumers reconstruct from observed
behavior. The two disciplined answers to ``who writes it first?'' define the
two workflows:

\begin{definition}[Design-first (contract-first)]
The OpenAPI document is authored \emph{before} implementation, reviewed as an
artifact in its own right, and treated as the source of truth. Server stubs,
client SDKs, contract tests, and mock servers are generated from it; the
implementation is continuously checked for conformance to it.
\end{definition}

\begin{definition}[Code-first (implementation-first)]
The contract is \emph{derived} from annotated code. FastAPI is the canonical
code-first toolchain: type annotations on path operations and Pydantic models
are compiled into an OpenAPI~3.1 document at runtime
(\texttt{GET /openapi.json})~\cite{fastapidocs}.
\end{definition}

Neither is morally superior; they optimize for different failure modes.
Design-first pays a front-loaded cost --- you must design representations
before you can test an idea --- but buys parallel workstreams (consumers build
against mocks from day one), independent contract review, and immunity to
``accidental contract'' drift where a refactor silently changes the wire
shape. Code-first pays almost nothing up front and keeps the contract
bit-identical to the implementation by construction, but risks shipping an
accidental contract: whatever the code happens to emit today becomes, de
facto, a promise. The pragmatic synthesis used throughout this book:
\emph{code-first mechanics, design-first governance}. Author the API in
FastAPI, but treat the generated \texttt{openapi.json} as a reviewed,
versioned, CI-asserted artifact --- the test suite in
Section~\ref{sec:ch6-code-tests} asserts its contents exactly.

\begin{table}[h]
  \centering
  \small
  \begin{tabularx}{\textwidth}{l X X}
    \toprule
    & \textbf{Design-first} & \textbf{Code-first (FastAPI)} \\
    \midrule
    Source of truth & Hand-written OpenAPI YAML & Type annotations + Pydantic models \\
    Contract fidelity & Drift possible (code may diverge) & Exact by construction \\
    Consumer start time & Immediately (generated mocks) & After first deployable build \\
    Review surface & Diff of the spec file & Diff of code \emph{and} pinned spec \\
    Best when & Public API, many external consumers & Internal APIs, small teams, fast iteration \\
    Main hazard & Spec/implementation skew & Accidental contract, unreviewed drift \\
    \bottomrule
  \end{tabularx}
  \caption{Design-first vs.\ code-first: the honest trade-off.}
  \label{tab:ch6-design-first}
\end{table}

\begin{figure}[h]
  \centering
  \begin{tikzpicture}[
      font=\small,
      box/.style={rectangle, rounded corners=2pt, draw=darkblue, fill=blue!5, minimum width=4.6cm, minimum height=0.8cm, align=center},
      art/.style={rectangle, rounded corners=2pt, draw=packtgray, fill=lightgray, minimum width=4.6cm, minimum height=0.7cm, align=center},
      flow/.style={->, >=stealth, shorten >=1pt},
      lbl/.style={font=\scriptsize\itshape, fill=white, inner sep=1pt}]
    % Design-first lane (top)
    \node[box] (spec) {author \texttt{openapi.yaml}};
    \node[box, right=1.1cm of spec] (gen) {generate stubs, mocks, SDKs};
    \node[box, right=1.1cm of gen] (implA) {implement against contract};
    \node[art, below=0.7cm of implA] (gateA) {CI: conformance gate};
    \draw[flow] (spec) -- (gen);
    \draw[flow] (gen) -- (implA);
    \draw[flow] (implA) -- (gateA);
    \node[note, left=0.25cm of spec] {\textbf{design-first}};

    % Code-first lane (bottom)
    \node[box, below=2.6cm of spec] (code) {annotated FastAPI code};
    \node[box, right=1.1cm of code] (derive) {derive \texttt{openapi.json}};
    \node[box, right=1.1cm of derive] (review) {pin, review, publish contract};
    \node[art, below=0.7cm of review] (gateB) {CI: schema assertions};
    \draw[flow] (code) -- (derive);
    \draw[flow] (derive) -- (review);
    \draw[flow] (review) -- (gateB);
    \node[note, left=0.25cm of code] {\textbf{code-first}};
  \end{tikzpicture}
  \caption{The two contract workflows converge on the same discipline: a
  reviewed OpenAPI artifact and a CI gate. They differ only in which side of
  the arrow is edited by hand.}
  \label{fig:ch6-workflows}
\end{figure}

\subsection{OpenAPI 3.1 as the Backbone}
\label{sec:ch6-openapi}

OpenAPI~3.1~\cite{openapi31} describes an HTTP API as a JSON or YAML document
with a small set of top-level keys. Know them cold:

\begin{itemize}
  \item \textbf{\texttt{openapi}} --- the version string (\texttt{"3.1.0"}).
        3.1's headline change: schemas are now \emph{full JSON Schema
        2020-12} (Chapter 4), replacing the 3.0 subset with its awkward
        \texttt{nullable} keyword.
  \item \textbf{\texttt{info}} --- \texttt{title}, \texttt{version}, contact
        and license metadata. The version here is \emph{your API's} version,
        not the spec's.
  \item \textbf{\texttt{servers}} --- an array of base URLs
        (\texttt{url} plus \texttt{description}, optionally with templated
        variables). This is how generated clients and the interactive docs
        know where the API lives in each environment.
  \item \textbf{\texttt{paths}} --- the heart: a map from path templates
        (\texttt{/orders/\{order\_id\}}) to \emph{path items}, each holding
        one \emph{operation} object per HTTP method. An operation carries
        \texttt{operationId}, \texttt{tags}, \texttt{parameters},
        \texttt{requestBody}, and --- most importantly --- \texttt{responses}:
        a map from status code to a representation schema. A complete
        operation documents its error shapes too; FastAPI emits the 422
        validation response automatically.
  \item \textbf{\texttt{components}} --- the reusable parts bin:
        \texttt{schemas} (Pydantic models land here, referenced by
        \texttt{\$ref}), \texttt{parameters}, \texttt{responses},
        \texttt{requestBodies}, \texttt{headers}, and
        \texttt{securitySchemes}.
  \item \textbf{\texttt{securitySchemes} / \texttt{security}} --- declarative
        descriptions of authentication mechanisms (\texttt{apiKey},
        \texttt{http} bearer/basic, \texttt{oauth2}, \texttt{openIdConnect})
        and where they apply. Chapter 16 implements enforcement; here we
        publish the scheme in the contract.
  \item \textbf{\texttt{tags}} --- named groups with descriptions; they drive
        the grouping in generated documentation.
\end{itemize}

FastAPI generates all of this from your annotations: the
\texttt{response\_model} becomes the success response schema, path and query
parameters become \texttt{parameters}, the request model becomes
\texttt{requestBody}, and every Pydantic model is hoisted into
\texttt{components.schemas}. The customization hook
(Section~\ref{sec:ch6-code-orders}) replaces \texttt{app.openapi} to inject
\texttt{servers}, a \texttt{securitySchemes} entry, and extensions such as
\texttt{x-audience} --- the escape hatch that keeps generated contracts
publishable as organizational artifacts.

\subsection{ASGI and the Server Stack}
\label{sec:ch6-asgi}

FastAPI code never touches a socket. Beneath it sits a three-layer stack, and
knowing exactly where each layer ends is the difference between debugging
production and guessing at it.

\begin{definition}[ASGI]
The \emph{Asynchronous Server Gateway Interface}: a callable contract between
an HTTP server and a Python application. An ASGI application is
\texttt{async def app(scope, receive, send)}: \texttt{scope} is a dict of
connection-lifetime metadata (method, path, headers, the \texttt{state} bag);
\texttt{receive} is an awaitable returning the next event
(\texttt{http.request} body chunks, \texttt{http.disconnect});
\texttt{send} awaits outbound events (\texttt{http.response.start} with
status and headers, then \texttt{http.response.body} chunks).
\end{definition}

\begin{enumerate}
  \item \textbf{The server} (uvicorn or hypercorn) owns TCP sockets, TLS
        termination, HTTP/1.1 parsing, and --- in uvicorn's
        \texttt{[standard]} build --- the \texttt{uvloop} event loop policy
        and \texttt{httptools} parser. Its whole job is translating bytes on
        the wire into ASGI events and back. Uvicorn is the default choice
        (C-accelerated, minimal); hypercorn adds HTTP/2 and HTTP/3 (QUIC)
        support.
  \item \textbf{Starlette} is the ASGI toolkit: the router that matches
        \texttt{scope["path"]}, the middleware machinery, request/response
        abstractions, \texttt{TestClient}, and the \emph{lifespan} protocol
        (startup/shutdown as ASGI messages).
  \item \textbf{FastAPI} is a routing subclass plus a dependency-injection
        and validation compiler on top of Starlette. When you write
        \texttt{@app.get("/orders")}, you are registering a Starlette route
        whose endpoint wrapper builds the DI graph, runs Pydantic validation,
        calls your function, and serializes the result.
\end{enumerate}

\begin{figure}[h]
  \centering
  \begin{tikzpicture}[
      font=\small,
      stage/.style={rectangle, rounded corners=2pt, draw=darkblue, fill=blue!5, minimum width=3.3cm, minimum height=0.9cm, align=center},
      ext/.style={rectangle, rounded corners=2pt, draw=packtgray, fill=lightgray, minimum width=2.2cm, minimum height=0.9cm, align=center},
      flow/.style={->, >=stealth, shorten >=1pt},
      lbl/.style={font=\scriptsize\itshape, fill=white, inner sep=1pt}]
    \node[ext] (client) {HTTP client};
    \node[stage, right=0.9cm of client] (uvicorn) {uvicorn\\[-2pt]{\scriptsize sockets, TLS, parsing}};
    \node[stage, right=0.9cm of uvicorn] (mw) {middleware stack\\[-2pt]{\scriptsize correlation, timing}};
    \node[stage, right=0.9cm of mw] (router) {Starlette router\\[-2pt]{\scriptsize path match}};
    \node[stage, below=1.2cm of router] (deps) {DI resolution\\[-2pt]{\scriptsize Depends graph}};
    \node[stage, left=0.9cm of deps] (valid) {Pydantic validation\\[-2pt]{\scriptsize 422 on failure}};
    \node[stage, left=0.9cm of valid] (handler) {handler\\[-2pt]{\scriptsize def $\to$ threadpool}};
    \node[stage, left=0.9cm of handler] (serial) {serialization\\[-2pt]{\scriptsize response\_model filter}};

    \draw[flow] (client) -- node[lbl, above]{bytes} (uvicorn);
    \draw[flow] (uvicorn) -- node[lbl, above]{scope / receive / send} (mw);
    \draw[flow] (mw) -- (router);
    \draw[flow] (router) -- (deps);
    \draw[flow] (deps) -- (valid);
    \draw[flow] (valid) -- (handler);
    \draw[flow] (handler) -- (serial);
    \draw[flow] (serial.south) to [bend right=30] node[lbl, below]{http.response.start / body} (client.south);
  \end{tikzpicture}
  \caption{The ASGI request lifecycle. Left of the router is the server; the
  router and middleware are Starlette; DI, validation, and response-model
  serialization are FastAPI.}
  \label{fig:ch6-asgi-pipeline}
\end{figure}

\textbf{Event-loop implications.} Uvicorn runs one OS thread per worker
process, and that thread runs one event loop. Every coroutine --- every
middleware, every \texttt{async def} endpoint, every streaming generator ---
shares it. A coroutine that blocks (a synchronous database driver, a
\texttt{time.sleep}, a CPU-heavy serialization) does not block one request;
it blocks \emph{every} request the worker is serving. The escape hatch is the
threadpool, discussed next.

\subsection{Sync vs.\ Async Endpoints: Threadpool Semantics}
\label{sec:ch6-sync-async}

FastAPI dispatches on the \emph{declaration} of your endpoint:

\begin{itemize}
  \item \texttt{async def} endpoints are awaited directly \textbf{on the
        event loop}. They are correct when every blocking operation in the
        body is itself awaitable (async database driver, \texttt{httpx}
        async client, \texttt{asyncio.sleep}).
  \item \texttt{def} endpoints are submitted to a threadpool
        (anyio's limiter, default capacity \textbf{40 threads}) and awaited
        from the loop. The loop stays responsive; the cost is a thread
        context switch and a hard ceiling of 40 concurrent blocking calls per
        worker.
\end{itemize}

This yields a decision rule with no exceptions: \emph{async when the I/O
stack is async end-to-end; def when any link in the chain blocks.} The
Orders API deliberately declares its CRUD endpoints \texttt{def}: the
\texttt{sqlite3} driver is synchronous, so calling it from
\texttt{async def} would stall the loop, while the threadpool absorbs it
correctly. The corollary is the capacity question: if p99 downstream latency
is $L$ ms and arrival rate is $\lambda$ requests/s, Little's Law says the
pool needs $\lambda \cdot L / 1000$ slots; 40 slots at 50 ms latency
saturates near 800 requests/s per worker --- beyond that, requests
\emph{queue inside the pool} and latencies climb without a single error.
Section~\ref{sec:ch6-case-study} is the war story of exactly that failure.

When does \texttt{async def} genuinely help? When concurrency per worker is
the bottleneck and the whole stack is awaitable: proxying, fan-out to several
upstream APIs (Chapter 2's async client patterns), WebSockets, and
server-streamed responses. For CPU-bound work neither helps --- that work
belongs in a subprocess or a task queue.

\textbf{Background tasks vs.\ task queues.} Starlette's
\texttt{BackgroundTasks} runs callables \emph{after} the response is sent, in
the same process: ideal for fire-and-forget side effects that may be lost on
deploy (cache warming, a best-effort audit write). It offers no retries, no
persistence, no visibility; a worker kill loses the task silently. Anything
whose loss is a business event --- sending the invoice, provisioning the
robot --- belongs in a real queue (Celery, arq, a broker-backed outbox;
Chapter 11 covers event-driven patterns). The decision is one question:
\emph{if this process dies a millisecond after the response, may the work be
lost?}

\subsection{The Dependency-Injection Machine}
\label{sec:ch6-di}

\texttt{Depends} is FastAPI's second compiler: from callable annotations it
builds, per request, a \emph{directed acyclic graph} of providers.

\begin{itemize}
  \item \textbf{Resolution.} A dependency is any callable --- function,
        class, or class instance --- whose own signature is analyzed for
        parameters, so dependencies compose into sub-dependencies. FastAPI
        resolves the graph depth-first, caching each provider's result
        \emph{per request}: if a handler and two of its sub-dependencies all
        request \texttt{get\_settings}, it executes \emph{once}. The DI
        showcase (Section~\ref{sec:ch6-code-di}) proves the caching
        empirically with a resolution counter.
  \item \textbf{Yield dependencies and cleanup.} A provider may
        \texttt{yield} its value; code after the \texttt{yield} runs after
        the response has been sent, in LIFO order, and runs \emph{even when
        the handler raised}. This is where connections close, sessions
        commit-or-rollback, and locks release. Pre-\texttt{yield} code runs
        during resolution; post-\texttt{yield} code is teardown --- never put
        response-affecting logic after the \texttt{yield}.
  \item \textbf{Scopes.} Dependencies are request-scoped by default
        (per-request caching). Process-lifetime singletons belong in
        \texttt{lifespan} state (as the Orders API does with its repository)
        or in \texttt{lru\_cache}-backed providers for stateless values such
        as settings.
  \item \textbf{Overrides.} \texttt{app.dependency\_overrides} maps the
        original provider to a replacement, keyed by the callable object
        itself. It is the canonical test seam: swap the repository for a
        stub, swap settings for test values, without touching route code ---
        demonstrated in the test suite (Section~\ref{sec:ch6-code-tests}).
\end{itemize}

\begin{figure}[h]
  \centering
  \begin{tikzpicture}[
      font=\small,
      dep/.style={rectangle, rounded corners=2pt, draw=darkblue, fill=blue!5, minimum width=3.4cm, minimum height=0.7cm, align=center},
      res/.style={rectangle, rounded corners=2pt, draw=packtgray, fill=lightgray, minimum width=3.4cm, minimum height=0.9cm, align=center},
      flow/.style={->, >=stealth, shorten >=1pt},
      lbl/.style={font=\scriptsize\itshape, fill=white, inner sep=1pt}]
    \node[dep] (route) {Starlette route match};
    \node[dep, below=0.9cm of route] (handler) {endpoint \texttt{create\_order}};
    \node[dep, below left=0.9cm and 0.3cm of handler] (repo) {\texttt{get\_repository}};
    \node[dep, below right=0.9cm and 0.3cm of handler] (settings) {\texttt{get\_settings}};
    \node[res, below=0.9cm of repo] (state) {\texttt{app.state.repository}\\[-2pt]{\scriptsize built once in lifespan}};
    \node[res, below=0.9cm of settings] (cfg) {\texttt{app.state.settings}\\[-2pt]{\scriptsize env-layered, 12-factor}};

    \draw[flow] (route) -- (handler);
    \draw[flow] (handler) -- node[lbl, above left]{Depends} (repo);
    \draw[flow] (handler) -- node[lbl, above right]{Depends} (settings);
    \draw[flow] (repo) -- (state);
    \draw[flow] (settings) -- (cfg);
  \end{tikzpicture}
  \caption{The Orders API dependency graph. Providers (white) resolve
  per-request and are cached once per request; resources (gray) are
  process-lifetime singletons constructed once in the lifespan and merely
  \emph{looked up} by providers. Overrides replace any provider node at test
  time.}
  \label{fig:ch6-di-graph}
\end{figure}

\subsection{Pydantic v2: The Validation Pipeline}
\label{sec:ch6-pydantic}

Pydantic v2~\cite{pydanticdocs} is a rewrite around \texttt{pydantic-core}, a
Rust validation engine; the Python layer is a thin schema builder. Two
consequences matter daily. First, \textbf{performance}: validation of a
realistic model is commonly 5--20$\times$ faster than v1 --- fast enough that
``validate everything at the boundary'' stops being a performance
conversation. Second, \textbf{the validator taxonomy}:

\begin{itemize}
  \item \texttt{@field\_validator("sku")} runs per field, before or after
        type coercion (\texttt{mode="before"} / \texttt{mode="after"}), and
        sees one field in isolation. Use it for per-field normalization and
        rules --- the Orders API lowercases \texttt{customer\_id} here.
  \item \texttt{@model\_validator(mode="after")} runs once the whole model is
        coerced, and sees \emph{all} fields. Use it for cross-field domain
        invariants --- the Orders API rejects duplicate SKU lines here,
        because no single field can see the duplication.
  \item \textbf{Serialization aliases.} \texttt{Field(serialization\_alias=)}
        decouples the Python name from the wire name; FastAPI serializes
        response models with \texttt{by\_alias=True}, so
        \texttt{customer\_id} becomes \texttt{customerId} on the wire while
        Python code keeps idiomatic names. \texttt{populate\_by\_name=True}
        keeps construction by Python name legal.
\end{itemize}

The final piece is the one with security consequences:
\texttt{response\_model}. FastAPI does not merely document it --- it
\emph{re-validates and filters} handler return values through it. Fields
absent from \texttt{response\_model} never reach the wire, which is why the
Orders API can return \texttt{OrderRecord} objects (carrying
\texttt{warehouse\_hint} and \texttt{fraud\_score}) from handlers and still
guarantee clients only ever see the \texttt{OrderRead} contract.

\subsection{Configuration, Logging, and the Operational Surface}
\label{sec:ch6-ops}

\textbf{12-factor configuration}~\cite{twelvefactor}: config that varies
between deploys lives in the environment, never in code.
\texttt{pydantic-settings} gives this a typed front door: a
\texttt{BaseSettings} subclass declares fields, defaults, and an
\texttt{env\_prefix} (\texttt{NWR\_}), and the library layers sources ---
init arguments, then environment variables, then a \texttt{.env} file, then
defaults --- with earlier sources winning. The layering is the point: tests
inject \texttt{Settings(repository\_backend="memory")} explicitly; production
sets \texttt{NWR\_SQLITE\_PATH}; nothing reads \texttt{os.environ} ad hoc.
\textbf{Secrets boundary:} settings objects may \emph{reference} secrets
(file paths, key names) but secret material itself enters through the
environment at deploy time and is never logged, never returned in a response,
and never committed --- the full treatment is Chapter 16.

\textbf{Structured logging with correlation IDs.} Text logs are for humans at
a laptop; production log pipelines ingest JSON. The Orders API formats every
record as one JSON object per line and binds a correlation ID --- the
incoming \texttt{X-Request-Id} or a minted one --- into a
\texttt{contextvars.ContextVar} at the outermost middleware, so every log
line emitted anywhere in the request, in any layer, carries the same ID.
Chapter 18 builds the full observability stack (metrics, tracing) on this
foundation.

\textbf{Probes.} Three distinct questions, three distinct endpoints:

\begin{table}[h]
  \centering
  \small
  \begin{tabularx}{\textwidth}{l X X}
    \toprule
    \textbf{Probe} & \textbf{Question} & \textbf{Checks} \\
    \midrule
    Startup & May the process begin receiving probes at all? & Lifespan completion: migrations, cache warmup, config load. \\
    Readiness (\texttt{/ready}) & May we route traffic to this instance? & Dependencies reachable (repository \texttt{ping()}), warmup done. \\
    Liveness (\texttt{/health}) & Is the process wedged beyond recovery? & Nothing external --- prove the loop runs; fail $\Rightarrow$ restart. \\
    \bottomrule
  \end{tabularx}
  \caption{Probe semantics. Conflating readiness with liveness causes boot
  loops: a dependency outage makes liveness fail, the orchestrator restarts
  the pod, which changes nothing.}
  \label{tab:ch6-probes}
\end{table}

\textbf{Graceful shutdown.} Uvicorn's lifespan protocol sends
\texttt{lifespan.shutdown} on SIGTERM; the \texttt{asynccontextmanager}
resumes after its \texttt{yield}. The discipline: flip readiness to
unavailable \emph{first} (so the load balancer drains the instance), then
finish in-flight requests, then close repositories. The Orders API's
\texttt{lifespan} implements exactly this ordering.

\subsection{Response Engineering}
\label{sec:ch6-responses}

The response side of the contract deserves the same rigor as the request
side:

\begin{itemize}
  \item \textbf{\texttt{response\_model} filtering} --- the security boundary
        from Section~\ref{sec:ch6-pydantic}: declare the outbound shape, let
        FastAPI drop everything else.
  \item \textbf{Status codes} --- 201 with the created representation for
        POST, 204 with an empty body for DELETE, 422 for validation failure,
        409 for domain-state conflicts (the invalid transition), 404 for
        missing resources, 503 for unready dependencies. RFC~9110~\S15
        defines the semantics~\cite{rfc9110}; Chapter 7 deepens error bodies
        into full RFC~9457 problem details.
  \item \textbf{Custom responses} --- return a \texttt{Response} subclass
        directly to bypass serialization (the error handlers return
        \texttt{JSONResponse} with \texttt{application/problem+json}), or
        subclass \texttt{JSONResponse} for alternative encoders.
  \item \textbf{Streaming} --- \texttt{StreamingResponse} with an iterator
        sends body chunks as generated, in constant memory; a \emph{sync}
        generator runs on the threadpool, an \texttt{async} generator on the
        loop. The NDJSON endpoint in Section~\ref{sec:ch6-code-orders}
        snapshots the records first, then streams one JSON object per line.
  \item \textbf{ETags and caching} --- validators and conditional requests
        are middleware-level concerns; Chapter 10 builds the ETag middleware
        this chapter points at.
\end{itemize}

\begin{examplebox}
\textbf{The accidental-contract trap, concretely.} Without
\texttt{response\_model}, returning an \texttt{OrderRecord} would serialize
\texttt{fraud\_score} --- an internal risk signal --- into every client
response. No exception is raised; the leak ships silently. Declaring
\texttt{response\_model=OrderRead} makes the leak structurally impossible,
and the test suite asserts absence of the field. This is why ``always declare
\texttt{response\_model}'' is a rule, not a style preference.
\end{examplebox}

%% ----------------------------------------------------------------------------
\section{Production-Ready Code Implementation}
\label{sec:ch6-code}

The chapter ships four complete, runnable Python files. Layout:

\begin{minted}{text}
orders_api.py        # the centerpiece: models, repositories, DI, routes
middleware.py        # pure-ASGI correlation-ID + timing middleware
di_showcase.py       # DI mechanics: chains, yield cleanup, caching, overrides
test_orders_api.py   # offline TestClient suite (17 tests)
\end{minted}

All four run offline; the test suite completes in about a second. Save them
in one directory, then (PowerShell):

\begin{minted}{bash}
python -m venv .venv
.\.venv\Scripts\Activate.ps1
pip install "fastapi>=0.110" "uvicorn[standard]>=0.29" "pydantic>=2.7" `
            "pydantic-settings>=2.2" "httpx>=0.27" "pytest>=8"
pytest test_orders_api.py -q
\end{minted}

\subsection{The Centerpiece: The Northwind Robotics Orders API}
\label{sec:ch6-code-orders}

Read the file top to bottom in the order it is written: settings, logging,
domain models, repositories, lifespan, dependencies, error mapping, then the
application factory. Each layer only depends on layers above it --- that
ordering discipline is what keeps the DI graph acyclic and the file readable
in one pass.

\begin{minted}{python}
r"""Northwind Robotics Orders API -- production-grade FastAPI reference service.

Everything runs offline and deterministically. Persistence is either an
in-memory repository (default) or SQLite, selected by settings.

Run locally (PowerShell):
    .\.venv\Scripts\Activate.ps1
    uvicorn orders_api:app --port 8000
Exercise it:
    curl.exe http://localhost:8000/health
    curl.exe -X POST http://localhost:8000/orders -H "Content-Type: application/json" -d '{"customer_id":"acme","lines":[{"sku":"NWR-1001","quantity":2,"unit_price_cents":4999}]}'
Test it fully in-process (no network):
    pytest test_orders_api.py -q
"""
from __future__ import annotations

import json
import logging
import sqlite3
import threading
from collections.abc import AsyncIterator, Iterator
from contextlib import asynccontextmanager
from datetime import datetime, timezone
from enum import StrEnum
from itertools import count
from typing import Annotated, Any, Protocol

from fastapi import Depends, FastAPI, Query, Request, Response, status
from fastapi.exceptions import RequestValidationError
from fastapi.openapi.utils import get_openapi
from fastapi.responses import JSONResponse, StreamingResponse
from pydantic import BaseModel, ConfigDict, Field, field_validator, model_validator
from pydantic_settings import BaseSettings, SettingsConfigDict

from middleware import CorrelationIdMiddleware, TimingMiddleware, request_id_var

logger = logging.getLogger("northwind.orders")

# ---------------------------------------------------------------------------
# Configuration (12-factor: everything environment-driven, nothing hardcoded)
# ---------------------------------------------------------------------------


class Settings(BaseSettings):
    """Environment-layered configuration; NWR_* env vars override defaults."""

    model_config = SettingsConfigDict(
        env_prefix="NWR_", env_file=".env", extra="ignore"
    )

    app_name: str = "northwind-orders-api"
    app_version: str = "1.0.0"
    environment: str = "dev"  # dev | staging | prod
    repository_backend: str = "memory"  # memory | sqlite
    sqlite_path: str = "orders.db"
    log_level: str = "INFO"


# ---------------------------------------------------------------------------
# Structured JSON logging (deepened in Chapter 18)
# ---------------------------------------------------------------------------


class JsonFormatter(logging.Formatter):
    """One JSON object per line, with the correlation ID of the request."""

    def format(self, record: logging.LogRecord) -> str:
        payload: dict[str, Any] = {
            "ts": datetime.fromtimestamp(record.created, tz=timezone.utc).isoformat(),
            "level": record.levelname,
            "logger": record.name,
            "msg": record.getMessage(),
            "request_id": request_id_var.get(),
        }
        for key in ("method", "path", "status_code", "elapsed_ms"):
            value = getattr(record, key, None)
            if value is not None:
                payload[key] = value
        if record.exc_info:
            payload["exc"] = self.formatException(record.exc_info)
        return json.dumps(payload, separators=(",", ":"), sort_keys=True)


def configure_logging(level: str) -> None:
    handler = logging.StreamHandler()
    handler.setFormatter(JsonFormatter())
    root = logging.getLogger()
    root.handlers[:] = [handler]
    root.setLevel(level.upper())


# ---------------------------------------------------------------------------
# Domain model (Pydantic v2)
# ---------------------------------------------------------------------------


class OrderStatus(StrEnum):
    PENDING = "pending"
    CONFIRMED = "confirmed"
    SHIPPED = "shipped"
    CANCELLED = "cancelled"


ALLOWED_TRANSITIONS: dict[OrderStatus, frozenset[OrderStatus]] = {
    OrderStatus.PENDING: frozenset({OrderStatus.CONFIRMED, OrderStatus.CANCELLED}),
    OrderStatus.CONFIRMED: frozenset({OrderStatus.SHIPPED, OrderStatus.CANCELLED}),
    OrderStatus.SHIPPED: frozenset(),  # terminal
    OrderStatus.CANCELLED: frozenset(),  # terminal
}


class OrderLine(BaseModel):
    """Immutable value object inside an order."""

    model_config = ConfigDict(frozen=True)

    sku: str = Field(pattern=r"^NWR-[0-9]{4}$", description="Northwind part SKU")
    quantity: int = Field(ge=1, le=1000)
    unit_price_cents: int = Field(ge=0, le=10_000_000)


class OrderCreate(BaseModel):
    """Inbound payload. Edge validation here; domain invariants in the model."""

    customer_id: str = Field(min_length=3, max_length=64)
    lines: list[OrderLine] = Field(min_length=1, max_length=50)

    @field_validator("customer_id")
    @classmethod
    def customer_id_must_be_a_slug(cls, value: str) -> str:
        normalized = value.strip().lower()
        if not normalized.replace("-", "").replace("_", "").isalnum():
            raise ValueError(
                "customer_id must be alphanumeric (dashes/underscores allowed)"
            )
        return normalized

    @model_validator(mode="after")
    def reject_duplicate_skus(self) -> "OrderCreate":
        skus = [line.sku for line in self.lines]
        if len(skus) != len(set(skus)):
            raise ValueError(
                "duplicate SKU lines are not allowed; merge quantities instead"
            )
        return self


class OrderStatusUpdate(BaseModel):
    status: OrderStatus


class OrderRead(BaseModel):
    """Outbound contract: exactly what clients are allowed to see."""

    model_config = ConfigDict(populate_by_name=True)

    id: str
    customer_id: str = Field(serialization_alias="customerId")
    lines: list[OrderLine]
    status: OrderStatus
    total_cents: int
    created_at: datetime
    version: int


class OrderRecord(OrderRead):
    """Stored record: a superset of the outbound contract.

    ``warehouse_hint`` and ``fraud_score`` are internal-only; they never
    leave the process because every route declares ``response_model=OrderRead``
    and the streaming endpoint excludes them explicitly.
    """

    warehouse_hint: str = "auto"
    fraud_score: float = 0.0


class OrderNotFoundError(LookupError):
    def __init__(self, order_id: str) -> None:
        super().__init__(f"order {order_id!r} not found")
        self.order_id = order_id


class InvalidTransitionError(ValueError):
    def __init__(self, current: OrderStatus, target: OrderStatus) -> None:
        super().__init__(f"cannot transition order from {current} to {target}")
        self.current = current
        self.target = target


def transition(record: OrderRecord, target: OrderStatus) -> OrderRecord:
    """Pure domain function: the ONLY place a status is allowed to change."""
    if target not in ALLOWED_TRANSITIONS[record.status]:
        raise InvalidTransitionError(record.status, target)
    return record.model_copy(
        update={"status": target, "version": record.version + 1}
    )


# ---------------------------------------------------------------------------
# Repository: the persistence boundary, injected as a dependency
# ---------------------------------------------------------------------------


class OrderRepository(Protocol):
    def new_id(self) -> str: ...
    def insert(self, record: OrderRecord) -> OrderRecord: ...
    def get(self, order_id: str) -> OrderRecord: ...
    def list(self, *, limit: int, offset: int) -> list[OrderRecord]: ...
    def replace(self, record: OrderRecord) -> OrderRecord: ...
    def delete(self, order_id: str) -> None: ...
    def ping(self) -> bool: ...
    def close(self) -> None: ...


class InMemoryOrderRepository:
    """Deterministic process-local store; IDs are sequential: ord-0001, ..."""

    def __init__(self) -> None:
        self._lock = threading.Lock()
        self._orders: dict[str, OrderRecord] = {}
        self._sequence = count(1)
        self._closed = False

    def new_id(self) -> str:
        with self._lock:
            return f"ord-{next(self._sequence):04d}"

    def insert(self, record: OrderRecord) -> OrderRecord:
        with self._lock:
            self._orders[record.id] = record
        return record

    def get(self, order_id: str) -> OrderRecord:
        try:
            return self._orders[order_id]
        except KeyError:
            raise OrderNotFoundError(order_id) from None

    def list(self, *, limit: int, offset: int) -> list[OrderRecord]:
        with self._lock:
            records = list(self._orders.values())  # insertion order = stable
        return records[offset : offset + limit]

    def replace(self, record: OrderRecord) -> OrderRecord:
        with self._lock:
            if record.id not in self._orders:
                raise OrderNotFoundError(record.id)
            self._orders[record.id] = record
        return record

    def delete(self, order_id: str) -> None:
        with self._lock:
            if self._orders.pop(order_id, None) is None:
                raise OrderNotFoundError(order_id)

    def ping(self) -> bool:
        return not self._closed

    def close(self) -> None:
        self._closed = True


class SQLiteOrderRepository:
    """SQLite-backed store. sqlite3 is a BLOCKING driver: call it only from
    ``def`` endpoints, which FastAPI runs on the threadpool (Section 6.3)."""

    def __init__(self, path: str) -> None:
        self._lock = threading.Lock()
        self._conn = sqlite3.connect(path, check_same_thread=False)
        self._conn.execute(
            "CREATE TABLE IF NOT EXISTS orders ("
            "id TEXT PRIMARY KEY, seq INTEGER NOT NULL, doc TEXT NOT NULL)"
        )
        self._conn.commit()

    def new_id(self) -> str:
        with self._lock:
            row = self._conn.execute(
                "SELECT COALESCE(MAX(seq), 0) + 1 FROM orders"
            ).fetchone()
        return f"ord-{row[0]:04d}"

    def insert(self, record: OrderRecord) -> OrderRecord:
        seq = int(record.id.rsplit("-", 1)[1])
        with self._lock:
            self._conn.execute(
                "INSERT INTO orders (id, seq, doc) VALUES (?, ?, ?)",
                (record.id, seq, record.model_dump_json()),
            )
            self._conn.commit()
        return record

    def _row_to_record(self, row: tuple[str, str], order_id: str) -> OrderRecord:
        if row is None:
            raise OrderNotFoundError(order_id)
        return OrderRecord.model_validate_json(row[1])

    def get(self, order_id: str) -> OrderRecord:
        row = self._conn.execute(
            "SELECT id, doc FROM orders WHERE id = ?", (order_id,)
        ).fetchone()
        return self._row_to_record(row, order_id)

    def list(self, *, limit: int, offset: int) -> list[OrderRecord]:
        rows = self._conn.execute(
            "SELECT doc FROM orders ORDER BY seq LIMIT ? OFFSET ?", (limit, offset)
        ).fetchall()
        return [OrderRecord.model_validate_json(row[0]) for row in rows]

    def replace(self, record: OrderRecord) -> OrderRecord:
        with self._lock:
            cursor = self._conn.execute(
                "UPDATE orders SET doc = ? WHERE id = ?",
                (record.model_dump_json(), record.id),
            )
            self._conn.commit()
        if cursor.rowcount == 0:
            raise OrderNotFoundError(record.id)
        return record

    def delete(self, order_id: str) -> None:
        with self._lock:
            cursor = self._conn.execute(
                "DELETE FROM orders WHERE id = ?", (order_id,)
            )
            self._conn.commit()
        if cursor.rowcount == 0:
            raise OrderNotFoundError(order_id)

    def ping(self) -> bool:
        try:
            self._conn.execute("SELECT 1")
            return True
        except sqlite3.Error:
            return False

    def close(self) -> None:
        self._conn.close()


# ---------------------------------------------------------------------------
# Lifespan: startup and graceful shutdown
# ---------------------------------------------------------------------------


@asynccontextmanager
async def lifespan(app: FastAPI) -> AsyncIterator[None]:
    settings: Settings = app.state.settings
    configure_logging(settings.log_level)
    if settings.repository_backend == "sqlite":
        repository: OrderRepository = SQLiteOrderRepository(settings.sqlite_path)
    else:
        repository = InMemoryOrderRepository()
    app.state.repository = repository
    app.state.ready = True  # startup work done; readiness probe may pass
    logger.info(
        "startup complete",
        extra={"backend": settings.repository_backend, "env": settings.environment},
    )
    try:
        yield  # the application serves requests here
    finally:
        app.state.ready = False  # stop passing readiness before teardown
        repository.close()
        logger.info("shutdown complete")


# ---------------------------------------------------------------------------
# Dependencies
# ---------------------------------------------------------------------------


def get_settings(request: Request) -> Settings:
    return request.app.state.settings


def get_repository(request: Request) -> OrderRepository:
    return request.app.state.repository


SettingsDep = Annotated[Settings, Depends(get_settings)]
RepoDep = Annotated[OrderRepository, Depends(get_repository)]


# ---------------------------------------------------------------------------
# Error mapping: domain errors -> RFC 9457-shaped problem documents (ch. 7)
# ---------------------------------------------------------------------------


def problem(
    status_code: int, title: str, detail: str, **extra: Any
) -> JSONResponse:
    body = {
        "type": "about:blank",
        "title": title,
        "status": status_code,
        "detail": detail,
        **extra,
    }
    return JSONResponse(
        status_code=status_code,
        content=body,
        media_type="application/problem+json",
    )


async def order_not_found_handler(
    request: Request, exc: OrderNotFoundError
) -> JSONResponse:
    return problem(404, "Order not found", str(exc), order_id=exc.order_id)


async def invalid_transition_handler(
    request: Request, exc: InvalidTransitionError
) -> JSONResponse:
    return problem(
        409,
        "Invalid state transition",
        str(exc),
        current=exc.current.value,
        target=exc.target.value,
    )


async def validation_error_handler(
    request: Request, exc: RequestValidationError
) -> JSONResponse:
    # Sanitize: never forward raw Pydantic ctx objects to clients.
    errors = [
        {
            "loc": ".".join(str(part) for part in error["loc"]),
            "msg": error["msg"],
            "type": error["type"],
        }
        for error in exc.errors()
    ]
    return problem(422, "Validation failed", "Request payload is invalid", errors=errors)


async def unhandled_error_handler(request: Request, exc: Exception) -> JSONResponse:
    # Log the full stack trace server-side; return NONE of it to the client.
    logger.error("unhandled error on %s", request.url.path, exc_info=exc)
    return problem(500, "Internal server error", "An unexpected error occurred")


# ---------------------------------------------------------------------------
# Application factory
# ---------------------------------------------------------------------------


def create_app(settings: Settings | None = None) -> FastAPI:
    settings = settings or Settings()
    app = FastAPI(
        title="Northwind Robotics Orders API",
        version=settings.app_version,
        lifespan=lifespan,
        docs_url="/docs",
    )
    app.state.settings = settings
    app.state.ready = False

    # Last added = outermost. Correlation binds the ID first so the timing
    # log line and all downstream logs carry it.
    app.add_middleware(TimingMiddleware)
    app.add_middleware(CorrelationIdMiddleware)

    app.add_exception_handler(OrderNotFoundError, order_not_found_handler)
    app.add_exception_handler(InvalidTransitionError, invalid_transition_handler)
    app.add_exception_handler(RequestValidationError, validation_error_handler)
    app.add_exception_handler(Exception, unhandled_error_handler)

    # -- operational endpoints ------------------------------------------------

    @app.get("/health", tags=["ops"])
    async def health() -> dict[str, str]:
        # Liveness: is the process alive and able to run a coroutine?
        return {"status": "ok"}

    @app.get("/ready", tags=["ops"])
    def readiness(request: Request, repo: RepoDep) -> JSONResponse:
        # Readiness: MAY we receive traffic? Checks dependencies, not the loop.
        if not request.app.state.ready or not repo.ping():
            return JSONResponse(
                status_code=status.HTTP_503_SERVICE_UNAVAILABLE,
                content={"status": "unavailable"},
            )
        return JSONResponse(content={"status": "ready"})

    # -- orders CRUD ----------------------------------------------------------

    @app.post(
        "/orders",
        response_model=OrderRead,
        status_code=status.HTTP_201_CREATED,
        tags=["orders"],
    )
    def create_order(payload: OrderCreate, repo: RepoDep) -> OrderRecord:
        record = OrderRecord(
            id=repo.new_id(),
            customer_id=payload.customer_id,
            lines=payload.lines,
            status=OrderStatus.PENDING,
            total_cents=sum(
                line.quantity * line.unit_price_cents for line in payload.lines
            ),
            created_at=datetime.now(timezone.utc),
            version=1,
        )
        return repo.insert(record)

    @app.get("/orders", response_model=list[OrderRead], tags=["orders"])
    def list_orders(
        repo: RepoDep,
        limit: Annotated[int, Query(ge=1, le=100)] = 20,
        offset: Annotated[int, Query(ge=0)] = 0,
    ) -> list[OrderRecord]:
        return repo.list(limit=limit, offset=offset)

    @app.get("/orders/{order_id}", response_model=OrderRead, tags=["orders"])
    def get_order(order_id: str, repo: RepoDep) -> OrderRecord:
        return repo.get(order_id)

    @app.patch("/orders/{order_id}", response_model=OrderRead, tags=["orders"])
    def update_order_status(
        order_id: str, payload: OrderStatusUpdate, repo: RepoDep
    ) -> OrderRecord:
        current = repo.get(order_id)
        return repo.replace(transition(current, payload.status))

    @app.delete("/orders/{order_id}", status_code=status.HTTP_204_NO_CONTENT,
                tags=["orders"])
    def delete_order(order_id: str, repo: RepoDep) -> Response:
        repo.delete(order_id)
        return Response(status_code=status.HTTP_204_NO_CONTENT)

    @app.get("/orders-stream", tags=["orders"])
    def stream_orders(repo: RepoDep) -> StreamingResponse:
        # Snapshot first; the sync generator below runs on the threadpool.
        records = repo.list(limit=10_000, offset=0)

        def ndjson() -> Iterator[bytes]:
            for record in records:
                # Internal-only fields must not cross the wire on ANY path.
                line = record.model_dump_json(
                    by_alias=True, exclude={"warehouse_hint", "fraud_score"}
                )
                yield line.encode("utf-8") + b"\n"

        return StreamingResponse(ndjson(), media_type="application/x-ndjson")

    # -- OpenAPI customization -------------------------------------------------

    def custom_openapi() -> dict[str, Any]:
        if app.openapi_schema:
            return app.openapi_schema
        schema = get_openapi(
            title="Northwind Robotics Orders API",
            version=settings.app_version,
            description=(
                "Fictional order management for the Northwind Robotics "
                "parts universe. Synthetic data only."
            ),
            routes=app.routes,
            servers=[
                {"url": "http://localhost:8000", "description": "local dev"},
                {"url": "https://api.northwind-robotics.example",
                 "description": "production (fictional)"},
            ],
        )
        # Document the API-key scheme used by the edge gateway. Enforcement
        # itself is Chapter 16's subject; here we only publish the contract.
        schema.setdefault("components", {}).setdefault("securitySchemes", {})[
            "ApiKeyAuth"
        ] = {"type": "apiKey", "in": "header", "name": "X-API-Key"}
        schema["info"]["x-audience"] = "internal"
        app.openapi_schema = schema
        return app.openapi_schema

    app.openapi = custom_openapi  # type: ignore[method-assign]
    return app


# Module-level instance for ``uvicorn orders_api:app``.
app = create_app()
\end{minted}


Walk the file once more with the mechanics of
Section~\ref{sec:ch6-mechanics} in mind; six design decisions deserve a
second look:

\begin{enumerate}
  \item \textbf{The factory, not the module global.} \texttt{create\_app}
        takes a \texttt{Settings} and returns a configured \texttt{FastAPI}.
        Tests get isolated applications with deterministic state; uvicorn
        still gets a module-level \texttt{app}. This is the single most
        important structural habit in the file.
  \item \textbf{\texttt{OrderCreate} vs.\ \texttt{OrderRead} vs.\
        \texttt{OrderRecord}.} Three models, three audiences: inbound
        validation (with a cross-field \texttt{model\_validator}), outbound
        contract (with a serialization alias), and the internal superset.
        The \texttt{response\_model=OrderRead} declaration on every route is
        the firewall between the second and third.
  \item \textbf{Repositories behind a \texttt{Protocol}.} Routes never name
        a storage engine; they name the protocol and receive whatever the
        lifespan constructed (memory or SQLite) --- or whatever a test
        injected via \texttt{dependency\_overrides}. Deterministic,
        sequential IDs (\texttt{ord-0001}, \texttt{ord-0002}, \ldots) keep
        every test assertion exact; a production repository would mint
        UUIDv7s at the same seam.
  \item \textbf{\texttt{def} endpoints, on purpose.} \texttt{sqlite3} is a
        blocking driver, so the CRUD handlers are plain \texttt{def} and run
        on the threadpool (Section~\ref{sec:ch6-sync-async}). Only
        \texttt{/health} --- which touches nothing blocking --- is
        \texttt{async def}.
  \item \textbf{Errors are mapped, not raised ad hoc.} Domain exceptions
        (\texttt{OrderNotFoundError}, \texttt{InvalidTransitionError}) are
        raised deep in the call stack and translated to RFC~9457-shaped
        problem documents by registered handlers; the catch-all handler logs
        the full traceback server-side and returns none of it.
  \item \textbf{The OpenAPI hook publishes the contract.}
        \texttt{custom\_openapi} adds \texttt{servers}, an
        \texttt{ApiKeyAuth} security scheme (enforcement is Chapter 16's
        subject), and an \texttt{x-audience} extension, then caches the
        schema on \texttt{app.openapi\_schema} so generation happens once.
\end{enumerate}

\begin{keyconceptbox}
\textbf{The lifespan is the composition root.} Everything process-lifetime
--- the settings object, the repository, the logging configuration --- is
constructed in exactly one place (\texttt{lifespan}/\texttt{create\_app}) and
\emph{looked up} everywhere else. Dependencies look up; they do not build.
When you find a dependency constructing its own database pool per request,
the composition root has been bypassed.
\end{keyconceptbox}

\subsection{Middleware Deep-Dive: Correlation IDs and Timing}
\label{sec:ch6-code-middleware}

Both middlewares are written against the raw ASGI protocol --- no FastAPI or
Starlette imports --- which is exactly the point: at this layer the framework
\emph{is} the protocol. Study how \texttt{send} is wrapped to intercept
\texttt{http.response.start} (the only moment headers can still be added),
and how the correlation middleware guards non-HTTP scopes (lifespan,
websocket) before touching anything.

\begin{minted}{python}
"""Pure-ASGI middleware: correlation IDs and request timing.

Both classes implement the ASGI callable protocol directly --
``__call__(scope, receive, send)`` -- so nothing here depends on FastAPI
or Starlette. Registration order matters: the LAST middleware added is the
OUTERMOST (runs first on the request path, last on the response path):

    app.add_middleware(TimingMiddleware)          # added 1st -> inner
    app.add_middleware(CorrelationIdMiddleware)   # added 2nd -> outer

Because CorrelationIdMiddleware is outermost, the request-id context
variable is already bound when TimingMiddleware logs, so every access-log
line carries the correlation ID.
"""
from __future__ import annotations

import logging
import time
import uuid
from contextvars import ContextVar
from typing import Any, Awaitable, Callable, MutableMapping

# The three callables that make up the ASGI protocol surface.
ASGIScope = MutableMapping[str, Any]
ASGIMessage = MutableMapping[str, Any]
ASGIReceive = Callable[[], Awaitable[ASGIMessage]]
ASGISend = Callable[[ASGIMessage], Awaitable[None]]
ASGIApp = Callable[[ASGIScope, ASGIReceive, ASGISend], Awaitable[None]]

# Bound once per request by CorrelationIdMiddleware; read by the JSON log
# formatter so every log line inside the request carries the same ID.
request_id_var: ContextVar[str] = ContextVar("request_id", default="-")

logger = logging.getLogger("northwind.access")


class CorrelationIdMiddleware:
    """Extract (or mint) X-Request-Id, bind it, and echo it on the response."""

    def __init__(self, app: ASGIApp) -> None:
        self.app = app

    async def __call__(
        self, scope: ASGIScope, receive: ASGIReceive, send: ASGISend
    ) -> None:
        if scope["type"] != "http":  # pass lifespan/websocket through untouched
            await self.app(scope, receive, send)
            return

        headers = {k.lower(): v for k, v in scope.get("headers", [])}
        request_id = (
            headers.get(b"x-request-id", b"").decode("latin-1")
            or uuid.uuid4().hex[:16]
        )
        token = request_id_var.set(request_id)
        scope["request_id"] = request_id  # visible to the app via request.scope

        async def send_with_request_id(message: ASGIMessage) -> None:
            if message["type"] == "http.response.start":
                message.setdefault("headers", [])
                message["headers"].append(
                    (b"x-request-id", request_id.encode("latin-1"))
                )
            await send(message)

        try:
            await self.app(scope, receive, send_with_request_id)
        finally:
            request_id_var.reset(token)  # never leak IDs across requests


class TimingMiddleware:
    """Measure downstream wall time, add X-Process-Time-Ms, log one JSON line."""

    def __init__(self, app: ASGIApp) -> None:
        self.app = app

    async def __call__(
        self, scope: ASGIScope, receive: ASGIReceive, send: ASGISend
    ) -> None:
        if scope["type"] != "http":
            await self.app(scope, receive, send)
            return

        started = time.perf_counter()
        status_code = 0

        async def send_with_timing(message: ASGIMessage) -> None:
            nonlocal status_code
            if message["type"] == "http.response.start":
                status_code = message["status"]
                elapsed_ms = (time.perf_counter() - started) * 1000.0
                message.setdefault("headers", [])
                message["headers"].append(
                    (b"x-process-time-ms", f"{elapsed_ms:.2f}".encode("ascii"))
                )
            await send(message)

        try:
            await self.app(scope, receive, send_with_timing)
        finally:
            # Logs even on unhandled exceptions (status_code stays 0).
            elapsed_ms = (time.perf_counter() - started) * 1000.0
            logger.info(
                "request completed",
                extra={
                    "method": scope["method"],
                    "path": scope["path"],
                    "status_code": status_code,
                    "elapsed_ms": round(elapsed_ms, 2),
                    "request_id": request_id_var.get(),
                },
            )
\end{minted}


\textbf{Ordering semantics.} \texttt{add\_middleware} wraps the application
like layers of an onion: the \emph{last} middleware added is the
\emph{outermost} --- it runs first on the request path and last on the
response path. The Orders API adds \texttt{TimingMiddleware} first and
\texttt{CorrelationIdMiddleware} second, so correlation wraps timing: the
context variable is already bound when the timing middleware emits its
access-log line, and that line carries \texttt{request\_id}. Reverse the two
\texttt{add\_middleware} calls and every access log would show
\texttt{request\_id="-"}. The \texttt{finally} in both classes matters for
the same reason: an unhandled exception downstream must still reset the
context variable and must still produce a log line (with
\texttt{status\_code} left at 0 --- a silent, greppable marker of ``never
reached response start'').

\begin{warningbox}
\textbf{\texttt{BaseHTTPMiddleware} vs.\ pure ASGI.} Starlette's
\texttt{BaseHTTPMiddleware} is convenient (request/response objects instead
of raw messages) but historically buffered streaming responses, broke
\texttt{request.context} propagation into background tasks, and added a
task-hop of overhead per request. Pure-ASGI middleware, as above, has none
of those costs and forces you to learn the protocol you are operating on.
Reach for \texttt{BaseHTTPMiddleware} only when the request/response API
genuinely simplifies the code.
\end{warningbox}

\subsection{The DI Showcase: Chains, Yield-Cleanup, Caching, Overrides}
\label{sec:ch6-code-di}

A minimal, self-checking application that isolates the four DI mechanics.
Run it directly --- \texttt{python di\_showcase.py} --- and it proves its own
claims with \texttt{TestClient} assertions.

\begin{minted}{python}
"""Dependency-injection showcase: chains, yield cleanup, caching, overrides.

Self-checking demo (offline, in-process):
    python di_showcase.py
"""
from __future__ import annotations

import sqlite3
from collections.abc import Iterator
from typing import Annotated, Protocol

from fastapi import Depends, FastAPI
from fastapi.testclient import TestClient


# --- a tiny domain: a quote-of-the-day service --------------------------------


class QuoteSource(Protocol):
    def quote(self) -> str: ...


class HandbookQuoteSource:
    """The 'real' implementation: reads the Northwind SRE handbook."""

    def quote(self) -> str:
        return "Robots do not panic. -- Northwind SRE handbook"


QUOTE_RESOLUTIONS = {"count": 0}  # observable side effect for the caching demo


def get_quote_source() -> HandbookQuoteSource:
    QUOTE_RESOLUTIONS["count"] += 1
    return HandbookQuoteSource()


QuoteDep = Annotated[QuoteSource, Depends(get_quote_source)]


def get_audit_trail(source: QuoteDep) -> list[str]:
    """Sub-dependency: depends on the SAME QuoteSource the handler requests."""
    return [f"audit: served quote of length {len(source.quote())}"]


# --- a yield dependency with guaranteed cleanup --------------------------------


def get_connection() -> Iterator[sqlite3.Connection]:
    """One in-memory SQLite connection per request, closed afterwards."""
    conn = sqlite3.connect(":memory:", check_same_thread=False)
    conn.execute("CREATE TABLE visits (path TEXT)")
    conn.execute("INSERT INTO visits (path) VALUES ('/demo')")
    try:
        yield conn  # the endpoint (and sub-dependencies) run here
    finally:
        conn.close()  # runs even when the endpoint raises


ConnectionDep = Annotated[sqlite3.Connection, Depends(get_connection)]


def create_app() -> FastAPI:
    app = FastAPI(title="di-showcase")

    @app.get("/demo")
    def demo(
        direct: QuoteDep,  # requested here AND inside get_audit_trail ...
        audit: Annotated[list[str], Depends(get_audit_trail)],
        conn: ConnectionDep,
    ) -> dict[str, object]:
        rows = conn.execute("SELECT COUNT(*) FROM visits").fetchone()[0]
        return {
            "quote": direct.quote(),
            "audit": audit,
            "visits": rows,
            "quote_resolutions": QUOTE_RESOLUTIONS["count"],
        }

    return app


class FallbackQuoteSource:
    """Test double injected via dependency_overrides."""

    def quote(self) -> str:
        return "synthetic fallback quote"


if __name__ == "__main__":
    # 1) Per-request caching: the dependency graph requests QuoteSource twice
    #    (directly + via the audit sub-dependency) yet resolves it ONCE.
    with TestClient(create_app()) as client:
        first = client.get("/demo").json()
        second = client.get("/demo").json()
    assert first["quote_resolutions"] == 1, first
    assert second["quote_resolutions"] == 2, second
    print("per-request caching OK:", first["quote_resolutions"],
          "->", second["quote_resolutions"])

    # 2) Overrides: swap the implementation without touching route code.
    app = create_app()
    app.dependency_overrides[get_quote_source] = FallbackQuoteSource
    with TestClient(app) as client:
        body = client.get("/demo").json()
    assert body["quote"] == "synthetic fallback quote", body
    print("override OK:", body["quote"])
\end{minted}


Three observations. First, the \texttt{demo} handler requests
\texttt{QuoteSource} directly \emph{and} transitively through
\texttt{get\_audit\_trail}, yet \texttt{quote\_resolutions} increments once
per request --- per-request caching in action. Second, \texttt{get\_connection}
is a generator dependency: the connection lives exactly one request and its
\texttt{finally} runs even when the handler raises; this is the pattern for
sessions, transactions, and locks. Third, \texttt{dependency\_overrides} is
keyed by the \emph{original callable}, not by name --- override
\texttt{get\_quote\_source} and every dependent (including
\texttt{get\_audit\_trail}) sees the replacement, because resolution happens
through the graph, not through imports.

\subsection{The Test Suite: Offline, In-Process, Deterministic}
\label{sec:ch6-code-tests}

\texttt{TestClient}~\cite{httpxdocs} runs the real ASGI application
in-process --- routing, middleware, DI, validation, serialization --- with no
socket and no network. Used as a context manager it also drives the
\emph{lifespan}, so startup and graceful shutdown are under test too, not
just the routes. The suite covers the four contract surfaces that matter:
happy-path CRUD, validation errors (422), domain errors (404/409), and the
generated OpenAPI document itself.

\begin{minted}{python}
"""TestClient suite for the Orders API: fully offline, fully deterministic."""
from __future__ import annotations

import json

import pytest
from fastapi.testclient import TestClient

import di_showcase
from orders_api import (
    InMemoryOrderRepository,
    OrderNotFoundError,
    Settings,
    create_app,
    get_repository,
)

VALID_ORDER = {
    "customer_id": "Acme-Robotics",
    "lines": [{"sku": "NWR-1001", "quantity": 2, "unit_price_cents": 4999}],
}


@pytest.fixture()
def client() -> TestClient:
    # The context manager drives the lifespan: startup runs, readiness flips,
    # and shutdown/close is exercised on exit.
    app = create_app(Settings(repository_backend="memory"))
    with TestClient(app) as test_client:
        yield test_client


# -- operational endpoints ----------------------------------------------------


def test_health_and_readiness(client: TestClient) -> None:
    assert client.get("/health").json() == {"status": "ok"}
    assert client.get("/ready").json() == {"status": "ready"}


def test_readiness_fails_closed(client: TestClient) -> None:
    client.app.state.ready = False  # simulate mid-shutdown / failed warmup
    response = client.get("/ready")
    assert response.status_code == 503
    assert response.json() == {"status": "unavailable"}


# -- CRUD happy paths ----------------------------------------------------------


def test_create_order_filters_internal_fields(client: TestClient) -> None:
    response = client.post("/orders", json=VALID_ORDER)
    assert response.status_code == 201
    body = response.json()
    assert body["id"] == "ord-0001"  # deterministic IDs
    assert body["status"] == "pending"
    assert body["total_cents"] == 9998
    assert body["customerId"] == "acme-robotics"  # alias + normalization
    assert "customer_id" not in body
    # response_model=OrderRead strips internal-only fields:
    assert "warehouse_hint" not in body
    assert "fraud_score" not in body


def test_get_and_list_orders(client: TestClient) -> None:
    client.post("/orders", json=VALID_ORDER)
    client.post(
        "/orders",
        json={
            "customer_id": "globex",
            "lines": [{"sku": "NWR-2002", "quantity": 1, "unit_price_cents": 100}],
        },
    )
    fetched = client.get("/orders/ord-0002")
    assert fetched.status_code == 200
    assert fetched.json()["customerId"] == "globex"

    page = client.get("/orders", params={"limit": 1, "offset": 1})
    assert [o["id"] for o in page.json()] == ["ord-0002"]


def test_status_transition_happy_path(client: TestClient) -> None:
    client.post("/orders", json=VALID_ORDER)
    confirmed = client.patch("/orders/ord-0001", json={"status": "confirmed"})
    assert confirmed.status_code == 200
    assert confirmed.json()["status"] == "confirmed"
    assert confirmed.json()["version"] == 2
    shipped = client.patch("/orders/ord-0001", json={"status": "shipped"})
    assert shipped.json()["status"] == "shipped"


def test_delete_order_then_404(client: TestClient) -> None:
    client.post("/orders", json=VALID_ORDER)
    assert client.delete("/orders/ord-0001").status_code == 204
    assert client.get("/orders/ord-0001").status_code == 404


# -- validation and domain errors ----------------------------------------------


def test_invalid_sku_is_a_422_with_sanitized_errors(client: TestClient) -> None:
    bad = {"customer_id": "acme", "lines": [
        {"sku": "not-a-sku", "quantity": 1, "unit_price_cents": 100}]}
    response = client.post("/orders", json=bad)
    assert response.status_code == 422
    body = response.json()
    assert response.headers["content-type"].startswith("application/problem+json")
    assert body["title"] == "Validation failed"
    assert body["errors"][0]["loc"] == "body.lines.0.sku"
    assert "ctx" not in body["errors"][0]  # no raw internals leak


def test_duplicate_skus_rejected_by_model_validator(client: TestClient) -> None:
    dup = {"customer_id": "acme", "lines": [
        {"sku": "NWR-1001", "quantity": 1, "unit_price_cents": 100},
        {"sku": "NWR-1001", "quantity": 3, "unit_price_cents": 100}]}
    response = client.post("/orders", json=dup)
    assert response.status_code == 422
    assert "duplicate SKU" in response.json()["errors"][0]["msg"]


def test_missing_order_is_a_404_problem(client: TestClient) -> None:
    response = client.get("/orders/ord-9999")
    assert response.status_code == 404
    body = response.json()
    assert body["title"] == "Order not found"
    assert body["order_id"] == "ord-9999"


def test_invalid_transition_is_a_409(client: TestClient) -> None:
    client.post("/orders", json=VALID_ORDER)
    response = client.patch("/orders/ord-0001", json={"status": "shipped"})
    assert response.status_code == 409
    body = response.json()
    assert body["current"] == "pending"
    assert body["target"] == "shipped"


# -- middleware -----------------------------------------------------------------


def test_correlation_id_echoed(client: TestClient) -> None:
    response = client.get("/health", headers={"X-Request-Id": "req-abc-123"})
    assert response.headers["x-request-id"] == "req-abc-123"


def test_correlation_id_minted_when_absent(client: TestClient) -> None:
    response = client.get("/health")
    assert len(response.headers["x-request-id"]) == 16
    assert "x-process-time-ms" in response.headers


# -- OpenAPI contract ------------------------------------------------------------


def test_openapi_schema_is_customized_and_complete(client: TestClient) -> None:
    schema = client.get("/openapi.json").json()
    assert schema["info"]["title"] == "Northwind Robotics Orders API"
    assert schema["info"]["x-audience"] == "internal"
    assert {p for p in schema["paths"]} >= {
        "/orders", "/orders/{order_id}", "/health", "/ready", "/orders-stream"}
    assert schema["components"]["securitySchemes"]["ApiKeyAuth"] == {
        "type": "apiKey", "in": "header", "name": "X-API-Key"}
    assert schema["servers"][0]["url"] == "http://localhost:8000"
    create_op = schema["paths"]["/orders"]["post"]
    assert "201" in create_op["responses"]  # status codes documented for free
    assert "422" in create_op["responses"]


# -- streaming --------------------------------------------------------------------


def test_stream_orders_as_ndjson(client: TestClient) -> None:
    client.post("/orders", json=VALID_ORDER)
    client.post("/orders", json={
        "customer_id": "globex",
        "lines": [{"sku": "NWR-2002", "quantity": 5, "unit_price_cents": 50}]})
    response = client.get("/orders-stream")
    assert response.headers["content-type"].startswith("application/x-ndjson")
    lines = [json.loads(line) for line in response.text.splitlines()]
    assert [line["id"] for line in lines] == ["ord-0001", "ord-0002"]
    assert all("fraud_score" not in line for line in lines)


# -- the SQLite backend -------------------------------------------------------------


def test_sqlite_repository_roundtrip(tmp_path) -> None:
    settings = Settings(
        repository_backend="sqlite", sqlite_path=str(tmp_path / "orders.db")
    )
    with TestClient(create_app(settings)) as client:
        assert client.post("/orders", json=VALID_ORDER).status_code == 201
        body = client.get("/orders/ord-0001").json()
        assert body["total_cents"] == 9998
        # lifecycle: shutdown closed the DB, readiness now fails closed
    assert client.app.state.ready is False


# -- dependency overrides --------------------------------------------------------


class EmptyRepository(InMemoryOrderRepository):
    def get(self, order_id: str):  # every lookup misses
        raise OrderNotFoundError(order_id)


def test_dependency_override_swaps_the_repository() -> None:
    app = create_app(Settings(repository_backend="memory"))
    app.dependency_overrides[get_repository] = EmptyRepository
    with TestClient(app) as client:
        assert client.get("/orders/ord-0001").status_code == 404
    app.dependency_overrides.clear()


def test_di_showcase_caching_and_override() -> None:
    app = di_showcase.create_app()
    with TestClient(app) as client:
        first = client.get("/demo").json()
    assert first["quote_resolutions"] == 1  # resolved once per request

    app = di_showcase.create_app()
    app.dependency_overrides[di_showcase.get_quote_source] = (
        di_showcase.FallbackQuoteSource)
    with TestClient(app) as client:
        assert client.get("/demo").json()["quote"] == "synthetic fallback quote"
\end{minted}


Note what the suite asserts and what it refuses to assert. It asserts the
\emph{contract}: status codes, problem-document shapes, the aliased wire name
\texttt{customerId}, the \emph{absence} of internal fields, the correlation
header, and the contents of \texttt{/openapi.json}. It does not assert log
timestamps, elapsed times, or minted request-ID values --- those are runtime
facts, not contract facts, and pinning them would make the suite
non-deterministic. Deterministic repository IDs (\texttt{ord-0001}) are what
make every other assertion exact.

\begin{examplebox}
\textbf{Two seams, two styles.} The suite demonstrates both test seams from
this chapter. The \texttt{client} fixture uses the \emph{factory seam} ---
\texttt{create\_app(Settings(repository\_backend="memory"))} --- which is the
right choice when the whole environment changes. The override tests use the
\emph{override seam} --- \texttt{app.dependency\_overrides[get\_repository] =
EmptyRepository} --- which is the right choice when one collaborator changes
and everything else should stay real. Choosing per test keeps suites fast
\emph{and} honest.
\end{examplebox}

%% ----------------------------------------------------------------------------
\section{Common Pitfalls \& Anti-Patterns}
\label{sec:ch6-pitfalls}

Each of these has caused a real incident or a real data leak. They are
ordered roughly by how expensive they are to discover late.

\subsection*{Blocking the event loop in \texttt{async def} handlers}
The single most damaging FastAPI mistake. A synchronous database driver, a
\texttt{requests} call, or a \texttt{time.sleep} inside \texttt{async def}
does not block one request --- it blocks the worker's entire event loop:
every in-flight request stalls, health checks time out, and the orchestrator
concludes the pod is dead. The rule is structural, not stylistic: if any call
in the handler chain blocks, the handler is \texttt{def}. Audit dependencies
too --- a blocking \emph{dependency} under an \texttt{async} endpoint blocks
exactly the same loop.

\subsection*{Leaking stack traces in 500s}
Starlette's debug mode renders beautiful tracebacks --- including local
variables, which means connection strings and customer data --- into the HTTP
response. Production runs with debug off \emph{and} a registered
\texttt{Exception} handler that logs the traceback server-side and returns a
generic problem document. ``It only happens on unexpected errors'' is exactly
when you cannot predict what the frames contain.

\subsection*{Mutable default dependencies}
\texttt{def endpoint(cache: dict = \{\})} and module-level dependency
instances shared across requests are the FastAPI flavor of Python's classic
mutable-default-argument bug, with concurrency on top: two requests mutating
one dict is a data race that survives your test suite and detonates under
load. Providers should construct fresh state per request, or delegate
shared state to a deliberately synchronized process-lifetime resource built
in the lifespan (like the repository's \texttt{threading.Lock}).

\subsection*{Validation only at the edge, with no domain invariants}
Field constraints (\texttt{ge=1}, regex patterns) catch malformed
\emph{shapes}; they cannot catch malformed \emph{meaning}. ``Pending orders
cannot ship'' and ``an order may not contain the same SKU twice'' are domain
invariants: the first lives in the \texttt{transition} function, the second
in a \texttt{model\_validator}. Systems that validate only shapes end up
enforcing invariants in the database's error messages --- which is to say,
not enforcing them at all in any way a client can act on.

\subsection*{\texttt{response\_model} drift}
Declaring \texttt{response\_model} once and never auditing it rots: new
internal fields get added to the stored model, someone forgets the model is
returned from a route, and the field leaks. Defenses: keep the outbound
model (\texttt{OrderRead}) a distinct class from the stored record
(\texttt{OrderRecord}), and assert field absence in tests --- the suite's
\texttt{"fraud\_score" not in body} line is a security control, not a style
test.

\subsection*{Catching broad \texttt{Exception} in handlers}
\texttt{except Exception: return JSONResponse(...)} inside a handler swallows
\emph{programming} errors (a typo'd attribute, a violated invariant) together
with \emph{operational} errors, converting bugs into plausible-looking 200s
or 500s with no traceback anywhere. Raise specific domain exceptions; let the
registered handlers map them; let everything else reach the catch-all, where
it is logged with full context. Broad catches belong at exactly one layer:
the outermost error handler.

\subsection*{Heavy work at import time / module scope}
Connecting to a database, reading files, or building a repository at module
import means: tests cannot import the module without side effects, every
uvicorn worker re-does the work on fork, and configuration is frozen before
anyone could inject it. Everything expensive belongs in
\texttt{create\_app}/\texttt{lifespan}. The module level should contain only
declarations --- and the single, side-effect-free
\texttt{app = create\_app()} line for the server's benefit.

%% ----------------------------------------------------------------------------
\section{Hands-On Production Lab \& Real-World Case Study}
\label{sec:ch6-lab}

\subsection{Lab: Stand Up the Orders API and Interrogate Its Contract}
\label{sec:ch6-lab-run}

\textbf{Objective.} Run the Orders API locally, exercise it with
\texttt{curl.exe} and \texttt{TestClient}, and inspect the generated OpenAPI
document --- offline end to end.

\textbf{Setup} (PowerShell):

\begin{minted}{bash}
mkdir orders-lab; cd orders-lab
# save orders_api.py, middleware.py, di_showcase.py, test_orders_api.py here
python -m venv .venv
.\.venv\Scripts\Activate.ps1
pip install "fastapi>=0.110" "uvicorn[standard]>=0.29" "pydantic>=2.7" `
            "pydantic-settings>=2.2" "httpx>=0.27" "pytest>=8"
\end{minted}

\textbf{Step 1 --- run the suite.} \texttt{pytest test\_orders\_api.py -q}
should report 17 passed in about a second, with no network access. Read one
JSON log line from the output and find its \texttt{request\_id} field.

\textbf{Step 2 --- run the server.}

\begin{minted}{bash}
uvicorn orders_api:app --port 8000
\end{minted}

\textbf{Step 3 --- exercise it} (second PowerShell window; note
\texttt{curl.exe}, not the \texttt{curl} alias):

\begin{minted}{bash}
curl.exe http://localhost:8000/health
curl.exe http://localhost:8000/ready
curl.exe -X POST http://localhost:8000/orders -H "Content-Type: application/json" `
  -H "X-Request-Id: lab-001" `
  -d '{"customer_id":"acme","lines":[{"sku":"NWR-1001","quantity":2,"unit_price_cents":4999}]}'
curl.exe http://localhost:8000/orders/ord-0001
curl.exe -X PATCH http://localhost:8000/orders/ord-0001 -H "Content-Type: application/json" `
  -d '{"status":"shipped"}'
\end{minted}

Observe: the PATCH returns a 409 problem document (\texttt{pending} cannot go
to \texttt{shipped}); every response echoes \texttt{x-request-id: lab-001};
and the response body contains \texttt{customerId} --- never
\texttt{fraud\_score}.

\textbf{Step 4 --- inspect the contract.}

\begin{minted}{bash}
curl.exe http://localhost:8000/openapi.json -o openapi.json
\end{minted}

Open \texttt{openapi.json} and locate: the \texttt{servers} array injected by
\texttt{custom\_openapi}; the \texttt{ApiKeyAuth} entry under
\texttt{components.securitySchemes}; the 422 response on
\texttt{POST /orders}; and the \texttt{OrderRead} schema --- confirm
\texttt{warehouse\_hint} and \texttt{fraud\_score} appear nowhere in it.
Then open \texttt{http://localhost:8000/docs} and drive the API from the
generated Swagger UI.

\textbf{Step 5 --- flip the backend.} Stop the server, set
\texttt{\$env:NWR\_REPOSITORY\_BACKEND="sqlite"} and
\texttt{\$env:NWR\_SQLITE\_PATH="lab.db"}, restart, and repeat Step 3. No
code changed; the environment did. That is the 12-factor property made
visible. Stop with \texttt{Ctrl+C} and watch the shutdown log line --- the
lifespan's post-\texttt{yield} code running.

\textbf{Acceptance criteria.} You are done when: the suite passes; the 409
transition error reproduces; a request without \texttt{X-Request-Id} gets a
minted 16-hex-digit ID back; and you can point at the exact line in
\texttt{openapi.json} where the alias \texttt{customerId} is declared.

\subsection{Case Study: The Threadpool Exhaustion Incident}
\label{sec:ch6-case-study}

\textbf{Context.} A robotics-parts fulfillment service --- one uvicorn worker
class per pod, four pods --- served a warehouse scanning fleet. Its hottest
endpoint, \texttt{POST /shipments}, was declared \texttt{async def} because
``async is faster,'' but its body called the synchronous PostgreSQL driver
(\texttt{psycopg2} without the async wrapper) three times per request.

\textbf{Symptom.} Latency dashboards showed p99 climbing from 40 ms to 30 s
every day at the 09:00 scanning shift-change, with \emph{zero} 5xx errors ---
then mass pod restarts as liveness probes (served by the same event loop)
began timing out. Restarts made it worse: warming pods re-entered a queue
that was already saturated.

\textbf{Diagnosis.} The blocking driver calls ran \emph{on the event loop},
so each of the four workers could make progress on exactly one request at a
time. Effective concurrency per pod was one, not the hundreds the autoscaler
assumed; Little's Law did the rest. The tell was in the metrics: event-loop
lag spiked \emph{before} latency did, and per-pod CPU sat near one core ---
four busy loops, zero parallelism.

\textbf{The botched first fix.} An engineer converted the endpoint to
\texttt{def}, moving the driver calls to anyio's threadpool --- correct in
principle. But the 40-thread default was sized for incidental blocking, not
for the whole workload: at shift change the pool filled, requests queued
inside it, and latencies still climbed --- now without loop lag to explain
them. The second-order lesson: the threadpool is a \emph{bounded resource
with a queue}, and the queue is invisible unless you instrument pool
occupancy.

\textbf{Resolution.} Three changes, in order of leverage: (1)~the endpoint
stayed \texttt{def} and the three driver calls were batched into one
round-trip, cutting thread-seconds per request by three; (2)~the anyio
thread limiter was raised to match the measured concurrency need, with pool
occupancy exported as a metric and alerted; (3)~the roadmap item --- an
async driver end to end, restoring \texttt{async def} honestly --- was
scheduled rather than assumed free. Post-incident p99 at shift change: 180 ms.

\textbf{Takeaways.} \texttt{async} is a property of your I/O stack, not your
function keyword; threadpools are capacity you must size and instrument; and
liveness probes sharing the event loop turn overload into restart loops ---
one more reason \texttt{/health} checks nothing external.

%% ----------------------------------------------------------------------------
\section{Self-Assessment Exercises \& Deep-Dive Questions}
\label{sec:ch6-self-assessment}

Work these before the interview vault; they are designed to be answered with
code, not prose.

\begin{enumerate}[label=\textbf{E\arabic*.}, leftmargin=1.2cm, itemsep=0.6em]
  \item \textbf{Header dependency.} Add a dependency
        \texttt{get\_api\_version} that reads an \texttt{X-API-Version}
        header, defaults to \texttt{"1"}, and is available to every route via
        a typed alias. Write a test asserting the default and the override.
  \item \textbf{Idempotent create.} Extend \texttt{POST /orders} to accept an
        \texttt{Idempotency-Key} header: a repeated key with an identical
        body returns the original 201 response; the same key with a different
        body returns 409. (Chapter 10 generalizes this.)
  \item \textbf{Blocking-proof the loop.} Add a middleware that measures
        event-loop lag (schedule \texttt{asyncio.sleep(0)} and time the
        delay) and emits it as a log field. Demonstrate the lag with a
        deliberately blocking endpoint, then fix the endpoint.
  \item \textbf{Scope audit.} Write a websocket-free pure-ASGI middleware
        that counts requests per path and exposes the counts on
        \texttt{/metrics} (Prometheus text format). Why must the counter live
        on \texttt{app.state} rather than as a middleware attribute?
        (Hint: middleware instances are per-application; be precise about
        which \emph{object} would be wrong and why.)
  \item \textbf{Validator choreography.} Move the duplicate-SKU rule from
        \texttt{OrderCreate}'s \texttt{model\_validator} to the repository's
        \texttt{insert}, raising a domain error mapped to 422. Which location
        is correct, and what does your answer imply about invariants that
        span \emph{multiple} aggregates (e.g.\ ``a customer may have at most
        five open orders'')?
  \item \textbf{Streaming discipline.} Modify \texttt{/orders-stream} to page
        the repository (\texttt{limit}/\texttt{offset} batches of 100) inside
        the generator instead of snapshotting first. What failure mode does
        snapshotting prevent during a long stream with concurrent writes?
  \item \textbf{Shutdown honesty.} Write a test that opens a streaming
        response, exits the \texttt{TestClient} context mid-stream, and
        asserts the repository's \texttt{close()} ran. What does uvicorn do
        differently from \texttt{TestClient} here, and where is that
        documented?
  \item \textbf{Contract pinning.} Add a test that diffs
        \texttt{/openapi.json} against a checked-in snapshot and fails on any
        change. Then make a deliberate breaking change (rename a response
        field) and watch it catch you. This is design-first governance on a
        code-first toolchain.
\end{enumerate}

%% ----------------------------------------------------------------------------
\section{Interview Q\&A Vault}
\label{sec:ch6-interview}

Thirty-six questions, ordered junior $\to$ staff. Answer aloud before
reading.

\begin{enumerate}[label=\textbf{Q\arabic*.}, leftmargin=1.3cm, itemsep=0.8em]

\item \textbf{What is ASGI, and how does it differ from WSGI?}\\
The Asynchronous Server Gateway Interface is the contract between an HTTP
server and a Python application: a single awaitable callable
\texttt{app(scope, receive, send)}. WSGI is a synchronous, two-argument
callable (\texttt{environ, start\_response}) that handles exactly one
request per call and cannot represent websockets, server push, or
long-lived connections. ASGI's event-channel design supports all of those
and lets one process interleave thousands of concurrent connections on a
single event loop.

\item \textbf{Explain the ASGI \texttt{scope}/\texttt{receive}/\texttt{send}
protocol.}\\
\texttt{scope} is a connection-lifetime dict: \texttt{type}
(\texttt{"http"}, \texttt{"websocket"}, \texttt{"lifespan"}), method, path,
headers, plus the \texttt{state} bag shared with the lifespan.
\texttt{receive()} is an awaitable yielding the next inbound event ---
\texttt{http.request} body chunks (with \texttt{more\_body}) or
\texttt{http.disconnect}. \texttt{send()} awaits outbound events: exactly one
\texttt{http.response.start} (status + headers) followed by
\texttt{http.response.body} chunks. Every header mutation must happen before
\texttt{response.start} is sent --- that is why middleware wraps
\texttt{send} and intercepts that one message.

\item \textbf{Where does FastAPI end and Starlette begin?}\\
Starlette owns the ASGI surface: routing, middleware, requests, responses,
lifespan, \texttt{TestClient}. FastAPI is a Starlette subclass plus a
compiler: it analyzes your type annotations to build the dependency graph,
run Pydantic validation, serialize through \texttt{response\_model}, and
emit OpenAPI. If a feature touches sockets, paths, or the lifespan, it is
Starlette; if it touches your function signature, it is FastAPI.

\item \textbf{What happens when a \texttt{def} endpoint runs?}\\
FastAPI detects the plain function and, instead of awaiting it, submits it
to anyio's threadpool and awaits the \emph{future} on the event loop. The
loop stays responsive; the handler runs concurrently with other requests on
a worker thread. The pool is bounded (default 40 tokens), so at most 40
blocking calls run concurrently per worker --- beyond that, requests queue
inside the pool.

\item \textbf{When should an endpoint be \texttt{async def}?}\\
When every blocking operation in its call chain --- including dependencies
--- is awaitable: an async database driver, an async \texttt{httpx} client,
\texttt{asyncio.sleep}. If any link blocks, the endpoint must be
\texttt{def}. The keyword describes your I/O stack, not your performance
ambitions.

\item \textbf{What is the default threadpool size for \texttt{def}
endpoints, and how do you change it?}\\
Forty threads per worker (anyio's default limiter). Raise it in the lifespan
via anyio's thread limiter API
(\texttt{anyio.to\_thread.current\_default\_thread\_limiter().total\_tokens
= N}) and treat the number as a capacity decision you instrument, not a
magic constant.

\item \textbf{What happens if you call a blocking driver inside
\texttt{async def}?}\\
The event loop stalls for the duration of the call. Every concurrent
request on that worker --- including liveness probes --- freezes, latencies
climb without errors, and orchestrators start killing healthy-looking pods.
It is the classic FastAPI production incident; the fix is \texttt{def}
(threadpool) or a genuinely async driver.

\item \textbf{Walk the full request lifecycle of \texttt{POST /orders}.}\\
Uvicorn parses bytes and builds the HTTP scope; the middleware stack runs
outermost-first (correlation, then timing); the Starlette router matches
\texttt{/orders} + POST; FastAPI resolves the DI graph (\texttt{get\_settings},
\texttt{get\_repository}, cached per request); Pydantic validates
\texttt{OrderCreate} (422 on failure, short-circuiting the handler); the
\texttt{def} handler runs on the threadpool; the returned
\texttt{OrderRecord} is filtered through \texttt{response\_model=OrderRead}
and serialized by alias; middleware wraps up on the way out (timing logs,
correlation echoes the ID); uvicorn writes bytes.

\item \textbf{How does FastAPI's DI resolve sub-dependencies?}\\
Recursively: a dependency's own signature is analyzed, so
\texttt{Depends} can appear inside providers. FastAPI flattens the resulting
DAG at route-registration time and, per request, executes it depth-first
with a per-request cache keyed by the provider callable --- a shared
sub-dependency requested twice runs once and both consumers see the same
instance.

\item \textbf{What is per-request caching of dependencies, and when does it
not apply?}\\
Within one request, each provider executes at most once; the result is
reused everywhere it is depended upon. It does \emph{not} apply across
requests (no cross-request cache), and you can opt out per use with
\texttt{Depends(provider, use\_cache=False)} --- needed, for example, when a
provider mutates state and two consumers genuinely need independent
instances.

\item \textbf{How do \texttt{yield} dependencies work, and when does cleanup
run?}\\
The provider runs to its \texttt{yield}, the yielded value is injected, the
handler and response complete, and then the code after \texttt{yield} runs
--- in LIFO order across the graph, and in a \texttt{finally} sense: it runs
even when the handler raised. Use it for closing connections, committing or
rolling back transactions, and releasing locks. Never put logic after the
\texttt{yield} that must affect the response --- the response is already
sent.

\item \textbf{How do \texttt{dependency\_overrides} work?}\\
\texttt{app.dependency\_overrides} maps the original provider callable to a
replacement callable with a compatible signature. Resolution consults the
map first, so overriding a provider replaces it everywhere in the graph,
including inside sub-dependencies. Clear the map after the test
(\texttt{app.dependency\_overrides.clear()}); leaking overrides between
tests is a classic source of order-dependent suites.

\item \textbf{Pydantic v1 vs.\ v2: what changed, and why is validation
faster?}\\
v2 moved the validation engine to \texttt{pydantic-core}, written in Rust,
making typical model validation 5--20$\times$ faster. The Python API was
modernized: \texttt{@field\_validator}/\texttt{@model\_validator} replace
\texttt{@validator}/\texttt{@root\_validator}, \texttt{model\_dump()}/
\texttt{model\_validate()} replace \texttt{.dict()}/\texttt{parse\_obj()},
\texttt{model\_config = ConfigDict(...)} replaces \texttt{class Config},
and v1 shims raise hard errors rather than silently changing behavior.

\item \textbf{\texttt{field\_validator} vs.\ \texttt{model\_validator} ---
when each?}\\
Field validators see one field, before or after type coercion; use them for
per-field normalization and rules (trim, lowercase, format checks). Model
validators see the coerced whole model; use them for cross-field invariants
(duplicate detection, ``end after start''). Putting a cross-field rule in a
field validator fails because sibling fields may not exist yet.

\item \textbf{What is \texttt{serialization\_alias} and when do you use
it?}\\
It decouples the Python attribute name from the wire name: Python code uses
\texttt{customer\_id}; the JSON body carries \texttt{customerId}. FastAPI
serializes response models with \texttt{by\_alias=True}, so the alias is
what clients see. Use it when the contract's naming convention (camelCase)
differs from Python's, without giving up either.

\item \textbf{Why does \texttt{response\_model} matter for security?}\\
Because it is a filter, not just documentation: FastAPI re-validates handler
return values against it and drops undeclared fields. Returning an internal
record that happens to carry \texttt{fraud\_score} or password hashes is
safe only because the declared outbound model omits them. Without
\texttt{response\_model}, whatever the ORM object grew last sprint ships to
clients --- a leak that raises no exception.

\item \textbf{What happens if you omit \texttt{response\_model}?}\\
FastAPI serializes whatever the handler returns (via its JSON encoder) and
documents the response as unconstrained. Three costs: internal fields can
leak, the OpenAPI contract loses its response schema (breaking client
generation), and serialization bugs surface at runtime instead of at
response-validation time.

\item \textbf{How are 422 errors produced, and how do you customize
them?}\\
When Pydantic validation of path, query, header, or body parameters fails,
FastAPI raises \texttt{RequestValidationError} before the handler runs, and
the default handler returns 422 with the error array. Register your own
handler to reshape it (the Orders API emits an RFC~9457-shaped problem with
a sanitized \texttt{errors} array); sanitize --- raw Pydantic errors can
carry internal context objects that do not belong on the wire.

\item \textbf{Liveness vs.\ readiness vs.\ startup probes --- distinguish
them.}\\
Startup: has initialization finished (gates the other two). Liveness: is the
process wedged --- fail it and the orchestrator restarts you; check nothing
external. Readiness: may traffic be routed here --- checks dependencies and
warmup; fail it and you are drained, not restarted. Conflating readiness
into liveness turns a database outage into a restart loop.

\item \textbf{How does graceful shutdown work in uvicorn?}\\
On SIGTERM, uvicorn stops accepting connections, sends the ASGI
\texttt{lifespan.shutdown} message, and waits (bounded) for in-flight
requests. Your lifespan's post-\texttt{yield} code runs: flip readiness off
first so the load balancer drains you, then close repositories and flush
logs. SIGKILL after the grace period takes what it takes --- anything
unflushed is lost.

\item \textbf{\texttt{BackgroundTasks} vs.\ a real task queue?}\\
\texttt{BackgroundTasks} runs callables after the response, in-process: no
retries, no persistence, lost on deploy. Suitable for best-effort side
effects (cache warming, metrics flushes). Work whose loss is a business
event --- billing, provisioning, notification --- needs a durable queue with
retries and dead-lettering (Celery, arq, outbox pattern). The deciding
question: may this work vanish if the process dies one millisecond after the
response?

\item \textbf{Design-first vs.\ code-first --- argue both sides.}\\
Design-first: the OpenAPI document is authored and reviewed before code;
consumers start against mocks immediately; the contract cannot drift
silently --- at the cost of front-loaded design and a spec/implementation
skew risk. Code-first: annotations generate the contract, so it is exact by
construction and iteration is cheap --- at the cost of an ``accidental
contract'' that changes whenever code does. The synthesis: code-first
mechanics, design-first governance --- pin and review the generated schema
in CI.

\item \textbf{Name the top-level keys of an OpenAPI 3.1 document.}\\
\texttt{openapi} (spec version), \texttt{info} (title, API version),
\texttt{servers}, \texttt{paths} (path items holding per-method operations),
\texttt{components} (\texttt{schemas}, \texttt{parameters},
\texttt{responses}, \texttt{securitySchemes}, \ldots), \texttt{security},
\texttt{tags}, and \texttt{externalDocs}~\cite{openapi31}.

\item \textbf{How does FastAPI generate the OpenAPI schema, and how do you
customize it?}\\
At request time for \texttt{/openapi.json}, \texttt{app.openapi()} walks the
registered routes and compiles annotations into the document, caching the
result on \texttt{app.openapi\_schema}. Customize by assigning a replacement
callable to \texttt{app.openapi}: call \texttt{get\_openapi(...)}, mutate
the dict (servers, security schemes, \texttt{x-*} extensions), cache, and
return it. Respect the cache or you regenerate on every docs request.

\item \textbf{What changed between OpenAPI 3.0 and 3.1?}\\
3.1 aligned \texttt{components.schemas} with full JSON Schema 2020-12: real
\texttt{type} arrays (\texttt{type: [string, "null"]}) replace the 3.0
\texttt{nullable} keyword; \texttt{prefixItems}, \texttt{unevaluatedProperties},
and \texttt{\$defs} become legal; and the \texttt{examples} arrays semantics
were clarified. For API authors the practical headline is that one schema
dialect now serves validation, documentation, and codegen.

\item \textbf{Explain 12-factor configuration and pydantic-settings' source
layering.}\\
Config that varies per deploy lives in the environment, not in
code~\cite{twelvefactor}. \texttt{BaseSettings} declares typed fields with
defaults and an \texttt{env\_prefix}; resolution layers sources with
precedence: explicit init arguments, environment variables, \texttt{.env}
file, then field defaults. Tests inject init arguments; production injects
the environment; nobody parses \texttt{os.environ} ad hoc.

\item \textbf{Where do secrets live, and what must never be logged?}\\
Secret material enters via the environment (or a secrets manager the
environment points to) at deploy time; never in source control, never in
\texttt{.env} files committed to the repo, never in log lines, exception
messages, or error responses. Settings objects may reference secrets (paths,
key names); the debug-mode 500 page and careless \texttt{repr} of settings
are the two classic leak paths.

\item \textbf{Why an application factory instead of module-level app
construction?}\\
A factory (\texttt{create\_app(settings)}) gives every test an isolated,
deterministically configured application and makes the composition root
explicit. Module-level construction freezes configuration at import time,
couples import to side effects, and makes parallel tests share state. Keep
exactly one module-level line --- \texttt{app = create\_app()} --- for
\texttt{uvicorn module:app}.

\item \textbf{How do you test FastAPI offline?}\\
With \texttt{TestClient} (or httpx's \texttt{ASGITransport}): the real ASGI
app runs in-process --- routing, middleware, DI, validation, serialization
--- over no socket. Use the client as a context manager so the lifespan
(startup/shutdown) runs; inject test state through the factory or
\texttt{dependency\_overrides}; and assert the contract --- status codes,
error shapes, header echoes, and the OpenAPI document itself.

\item \textbf{How would you put a correlation ID on every log line of a
request?}\\
Outermost middleware extracts or mints \texttt{X-Request-Id} and binds it to
a \texttt{contextvars.ContextVar}; a logging \texttt{Formatter} reads the
variable for every record; the middleware echoes the ID on the response and
resets the variable in a \texttt{finally}. Context variables (not
thread-locals) because the request is a coroutine, not a thread.

\item \textbf{Explain middleware ordering semantics.}\\
\texttt{add\_middleware} wraps the app like an onion: the \emph{last} added
is \emph{outermost} --- first on the request path, last on the response
path. Order matters for causality: correlation must wrap timing if the
timing log should carry the request ID; authentication must wrap anything
that trusts identity.

\item \textbf{\texttt{BaseHTTPMiddleware} vs.\ pure-ASGI middleware?}\\
\texttt{BaseHTTPMiddleware} trades the raw protocol for request/response
objects --- convenient, but it historically buffered streaming responses,
added a task hop, and had context-propagation bugs with background tasks.
Pure-ASGI middleware (implement \texttt{\_\_call\_\_(scope, receive, send)})
has none of those costs and composes at exactly one layer. Prefer pure ASGI
unless the request/response API materially simplifies the code.

\item \textbf{Why is catching broad \texttt{Exception} inside handlers an
anti-pattern?}\\
It collapses programming errors and operational errors into one path: a
typo returns a plausible 500 (or worse, a 200-with-fallback) with the
traceback swallowed, and the error budget lies. Raise specific domain
exceptions, register handlers that map them to precise statuses, and let a
single catch-all log-and-sanitize everything unexpected. Broad catch at
exactly one layer.

\item \textbf{How does \texttt{StreamingResponse} work, and what are its
constraints?}\\
It takes an iterator of byte chunks and sends one
\texttt{http.response.body} message per chunk with \texttt{more\_body=True},
keeping memory constant regardless of payload size. A sync generator runs on
the threadpool; an \texttt{async} generator on the loop --- so an async
generator must not block. Headers (including \texttt{Content-Type}) are
committed at \texttt{response.start}, before the first chunk: you cannot
change your mind mid-stream, and mid-stream errors can only truncate, not
re-status.

\item \textbf{How would you size the threadpool for a def-heavy service?}\\
With Little's Law: required slots $\approx \lambda \cdot L$, with $\lambda$
the peak arrival rate and $L$ the mean blocking-service time per request in
seconds. Validate empirically: instrument pool occupancy and queue depth,
load-test at 2$\times$ expected peak, and set a limit that saturates
\emph{downstream} before it saturates memory --- a 400-thread pool in front
of a 50-connection database just moves the queue.

\item \textbf{How would you migrate a blocking production service to async
safely?}\\
Instrument first: event-loop lag, threadpool occupancy, downstream latency.
Migrate the leaves, not the trunk: replace the driver with its async
equivalent behind the repository protocol, keeping \texttt{def} endpoints.
Only when the whole chain is awaitable do you flip handlers to
\texttt{async def}, one route at a time, behind a config flag, with the
blocking variant still deployable. Gate the rollout on loop-lag SLOs; the
anti-pattern is the big-bang rewrite where a single missed blocking call
takes down the loop.

\end{enumerate}

%% ----------------------------------------------------------------------------
\section{External Bibliography \& Further Reading}
\label{sec:ch6-bibliography}

\begin{itemize}[itemsep=0.4em]
  \item \textbf{FastAPI documentation}~\cite{fastapidocs} --- the reference
        for path operations, dependencies, middleware, testing, and OpenAPI
        customization. Read the \emph{Advanced User Guide} sections on
        dependencies with \texttt{yield} and lifespan; they are this
        chapter's mechanics in the official voice.
  \item \textbf{Starlette documentation} --- the ASGI toolkit underneath:
        routing, middleware, \texttt{TestClient}, lifespan, and the
        pure-ASGI middleware style used in
        Section~\ref{sec:ch6-code-middleware}.
  \item \textbf{Uvicorn documentation} --- deployment, workers, the
        \texttt{[standard]} extras (\texttt{uvloop}, \texttt{httptools}),
        and the lifespan/timeout settings that govern graceful shutdown.
  \item \textbf{Pydantic v2 documentation}~\cite{pydanticdocs} ---
        validators, serialization aliases, \texttt{ConfigDict}, and the
        migration guide from v1. The \emph{pydantic-core} architecture notes
        explain where the 5--20$\times$ speedup comes from.
  \item \textbf{pydantic-settings documentation} --- source layering,
        \texttt{env\_prefix}, \texttt{.env} files, and secrets handling.
  \item \textbf{The OpenAPI 3.1 specification}~\cite{openapi31} --- the
        normative contract grammar; read the \emph{Paths}, \emph{Operation},
        \emph{Components}, and \emph{Security Scheme} objects alongside
        Section~\ref{sec:ch6-openapi}.
  \item \textbf{The Twelve-Factor App}~\cite{twelvefactor} --- factor III
        (config in the environment) is the configuration section's backbone;
        factor IX (disposability) is the graceful-shutdown section's.
  \item \textbf{RFC 9110 (HTTP Semantics)}~\cite{rfc9110} --- \S15 for the
        status-code semantics this chapter's routes implement; the status
        registry entries for 201, 204, 404, 409, and 422.
  \item \textbf{\emph{Building Microservices}, 2nd ed., Sam
        Newman}~\cite{newman2021microservices} --- the organizational frame:
        service boundaries, contract thinking, and why ``the API is the
        product'' outlives any framework choice.
  \item \textbf{httpx documentation}~\cite{httpxdocs} --- the
        \texttt{ASGITransport} mechanism underneath \texttt{TestClient}, and
        the async client patterns that Chapter 2 introduced.
\end{itemize}

%% ----------------------------------------------------------------------------
\section{Capstone Projects}
\label{sec:ch6-capstones}

Exactly five projects, progressive. All are offline, deterministic, and use
only synthetic Northwind Robotics data. Place deliverables under
\texttt{capstones\textbackslash{}capstone\_NN\_<slug>\textbackslash{}} per the
workspace conventions.

\subsection{Capstone 1: Todo API with a Full Test Suite [junior]}
\textbf{Objective.} Build a minimal but \emph{complete} FastAPI service: a
todo-list API with create, read, list, complete, and delete operations, an
in-memory repository behind a protocol, and a TestClient suite that covers
every route.\\
\textbf{Inputs/Datasets.} A seeded fixture of 20 synthetic todos (titles from
the Northwind Robotics maintenance log, e.g.\ ``recalibrate arm 7'').\\
\textbf{Steps.} (1)~Model \texttt{TodoCreate}/\texttt{TodoRead} with
constraints (title 3--80 chars; \texttt{priority} 1--5). (2)~Implement the
repository protocol and an in-memory implementation with deterministic IDs.
(3)~Wire routes with \texttt{response\_model} on all of them and correct
status codes (201/204/404/422). (4)~Write at least 10 tests: every happy
path, a 404, and two distinct 422s.\\
\textbf{Acceptance Criteria.}
\begin{itemize}
  \item[$\square$] \texttt{pytest -q} passes offline; IDs are deterministic
        per fresh app.
  \item[$\square$] Every route declares \texttt{response\_model}; no internal
        field appears in any response body.
  \item[$\square$] The OpenAPI schema documents a 422 response for every
        route that accepts a body.
\end{itemize}
\textbf{Stretch Goals.} Add the correlation-ID middleware and assert the
header echo; add a \texttt{/todos-stream} NDJSON endpoint with a streaming
test.

\subsection{Capstone 2: Inventory API with a SQLite Repository [mid]}
\textbf{Objective.} Rebuild the Orders pattern for warehouse inventory:
parts, stock levels, and reservations --- with \emph{both} repositories
(in-memory and SQLite) behind one protocol, selected by
\texttt{pydantic-settings}, and a readiness probe that actually exercises the
database.\\
\textbf{Inputs/Datasets.} A generated fixture of 500 synthetic parts
(\texttt{NWR-0001}\ldots\texttt{NWR-0500}) with stock counts; 1\,\%
deliberately invalid rows for negative tests.\\
\textbf{Steps.} (1)~Model \texttt{Part}, \texttt{Reservation}, and their
create/read counterparts. (2)~Implement both repositories; the SQLite one
must create its schema idempotently. (3)~Implement reservation as a
\emph{domain transition} with invariants (cannot reserve more than on-hand;
cannot cancel a shipped reservation) mapped to 409. (4)~Make \texttt{/ready}
fail when the database file is unreadable. (5)~Test both backends through
the factory seam.\\
\textbf{Acceptance Criteria.}
\begin{itemize}
  \item[$\square$] The same suite passes against \texttt{memory} and
        \texttt{sqlite} settings with zero code changes.
  \item[$\square$] Oversubscribing a part returns 409 with a problem document
        naming the part and the shortfall.
  \item[$\square$] Readiness returns 503 when the SQLite path is a directory
        (probing the failure path, not mocking it).
\end{itemize}
\textbf{Stretch Goals.} Add optimistic concurrency: a \texttt{version} field
checked on update (412 Precondition Failed on mismatch; Chapter 10
generalizes).

\subsection{Capstone 3: Plugin-Style Dependency-Injection Architecture [senior]}
\textbf{Objective.} Refactor the Orders API so that \emph{repositories,
notifiers, and pricing policies are plugins}: implementations register
themselves in a module-level registry, settings select them by name, and the
lifespan composes the chosen set --- with a test suite that swaps each plugin
via \texttt{dependency\_overrides} without touching route code.\\
\textbf{Inputs/Datasets.} The chapter's Orders API; two synthetic notifier
plugins (``webhook-log'' and ``audit-ndjson'') and two pricing policies
(``list'' and ``fleet-discount'').\\
\textbf{Steps.} (1)~Define protocols for \texttt{Notifier} and
\texttt{PricingPolicy}. (2)~Build a registry (name $\to$ factory) with
fail-fast errors on unknown names. (3)~Extend \texttt{Settings} with
\texttt{notifier: str} and \texttt{pricing: str}; the lifespan instantiates
the selection. (4)~Thread both through DI into the order lifecycle (price on
create, notify on status change). (5)~Prove isolation: each plugin has tests
using overrides, and one test boots every registry entry.\\
\textbf{Acceptance Criteria.}
\begin{itemize}
  \item[$\square$] Adding a new plugin touches exactly one new file plus one
        registry line --- no route or factory edits.
  \item[$\square$] Booting with \texttt{NWR\_NOTIFIER=bogus} fails at startup
        with an actionable error, not at first request.
  \item[$\square$] Override tests demonstrate that notifications and pricing
        are replaceable per test, per request graph.
\end{itemize}
\textbf{Stretch Goals.} Emit the effective plugin set and versions into
\texttt{/health} and into a startup log line; add a contract test that every
registered notifier satisfies the protocol against a shared behavioral suite.

\subsection{Capstone 4: Async Migration of a Blocking Service [senior]}
\textbf{Objective.} Take a deliberately blocking service (synchronous SQLite
repository, \texttt{time.sleep}-simulated latency) and migrate it to
correct-concurrency form, \emph{proving} each step with an event-loop-lag
middleware and a concurrency benchmark --- the Section~\ref{sec:ch6-case-study}
incident, reproduced and fixed in miniature.\\
\textbf{Inputs/Datasets.} A provided baseline service where
\texttt{async def} handlers call blocking code (the bug); a seeded dataset of
1{,}000 orders; a benchmark script issuing 200 concurrent requests through
\texttt{ASGITransport}.\\
\textbf{Steps.} (1)~Write the loop-lag middleware and the benchmark; capture
baseline p50/p99 and loop lag. (2)~Fix variant A: handlers become
\texttt{def} (threadpool). Re-measure. (3)~Fix variant B: replace the
repository with an \texttt{asyncio.to\_thread}-wrapped implementation behind
the same protocol and restore \texttt{async def}. Re-measure. (4)~Write the
decision memo: which variant ships, and what limiter setting, with the
numbers.\\
\textbf{Acceptance Criteria.}
\begin{itemize}
  \item[$\square$] Baseline shows loop lag $> 5\,$s under load; both fixed
        variants show lag $< 50\,$ms.
  \item[$\square$] p99 under 200-way concurrency improves by an order of
        magnitude; results are reproducible across three runs.
  \item[$\square$] The memo cites pool occupancy evidence, not folklore.
\end{itemize}
\textbf{Stretch Goals.} Add a third variant using a genuinely async store
(an \texttt{asyncio.Queue}-backed fake) and explain the remaining
differences; instrument and graph threadpool queue depth during the
transition.

\subsection{Capstone 5: OpenAPI-First Generator Round-Trip [staff]}
\textbf{Objective.} Close the design-first loop: author
\texttt{orders.openapi.yaml} by hand for a superset of the chapter's API
(adding shipments), generate a FastAPI skeleton and a typed Python client
from it, implement the skeleton, and build a CI-style gate that fails when
the running app's schema diverges from the authored contract --- then perform
a governed contract change end to end.\\
\textbf{Inputs/Datasets.} The authored spec; the chapter's Orders API as the
seed implementation; a ``change request'' adding an optional
\texttt{priority} field to orders (backward compatible) and one renaming a
field (breaking).\\
\textbf{Steps.} (1)~Author the 3.1 spec with \texttt{servers},
\texttt{components}, \texttt{securitySchemes}, and complete error responses.
(2)~Generate or hand-synthesize the skeleton from the spec; implement against
it. (3)~Implement the diff gate: fetch \texttt{/openapi.json}, compare
against the pinned spec, classify every difference as breaking or
compatible. (4)~Apply the compatible change: spec first, code second, gate
green. (5)~Attempt the breaking change and document how the gate stops it.\\
\textbf{Acceptance Criteria.}
\begin{itemize}
  \item[$\square$] The gate classifies the seeded compatible and breaking
        changes correctly, with human-readable reasons.
  \item[$\square$] The running service passes its own gate; the pinned spec
        and the generated schema round-trip with zero diffs.
  \item[$\square$] The typed client consumes the service in a TestClient
        integration test with no hand-edited models.
  \item[$\square$] A written policy states who may bump which version number
        and when (feeding Chapter 9 and Chapter 20).
\end{itemize}
\textbf{Stretch Goals.} Emit the diff report as an RFC~9457 problem document;
extend the gate to detect \emph{undocumented} endpoints (in code, absent
from spec) as a distinct violation class.

%% ----------------------------------------------------------------------------
\section{Python Dependencies Appendix}
\label{sec:ch6-deps}

Everything in this chapter pins to a minimal \texttt{requirements.txt}:

\begin{minted}{text}
fastapi>=0.110
uvicorn[standard]>=0.29
pydantic>=2.7
pydantic-settings>=2.2
httpx>=0.27
pytest>=8
\end{minted}

Install on Windows/PowerShell:

\begin{minted}{bash}
python -m venv .venv
.\.venv\Scripts\Activate.ps1
pip install -r requirements.txt
\end{minted}

\subsection{fastapi}
\textbf{Description.} The framework itself: annotation-driven routing,
dependency injection, Pydantic integration, and OpenAPI generation on top of
Starlette~\cite{fastapidocs}.\\
\textbf{Why included.} The subject of the chapter; every listing runs on it.\\
\textbf{Alternatives.} \textbf{Flask} (synchronous WSGI, vast ecosystem, no
native validation or contract generation), \textbf{Django Ninja}
(FastAPI-style typed APIs inside the Django universe --- choose when you
already pay for Django), \textbf{Litestar} (comparable typed-ASGI framework
with a batteries-included posture), \textbf{aiohttp} (lower-level async
framework for when you want to own the stack).\\
\textbf{Installation.} \texttt{pip install "fastapi>=0.110"}.

\subsection{uvicorn[standard]}
\textbf{Description.} The ASGI server: sockets, TLS, HTTP parsing, and the
event loop. The \texttt{[standard]} extra pulls in \texttt{uvloop} and
\texttt{httptools} --- the C-accelerated loop policy and parser used in
production.\\
\textbf{Why included.} Runs the app in the lab; its lifespan and shutdown
semantics are part of the chapter's operational surface. Not needed by the
test suite --- \texttt{TestClient} speaks ASGI directly, a fact worth
noticing.\\
\textbf{Alternatives.} \textbf{hypercorn} (HTTP/2 and HTTP/3 support),
\textbf{daphne} (Channels ecosystem), \textbf{granian} (Rust-based server).\\
\textbf{Installation.} \texttt{pip install "uvicorn[standard]>=0.29"} (quote
the brackets in PowerShell).

\subsection{pydantic}
\textbf{Description.} Data validation and settings-independent modeling:
the \texttt{BaseModel}, validators, and serialization machinery the whole
contract is built from~\cite{pydanticdocs}. v2's \texttt{pydantic-core} is a
Rust engine; validation speed stops being a design constraint.\\
\textbf{Why included.} Every inbound payload, outbound contract, and stored
record in the chapter is a Pydantic model; \texttt{response\_model} filtering
is a Pydantic behavior.\\
\textbf{Alternatives.} \textbf{attrs + cattrs} (class authoring plus
explicit (de)serialization; more assembly, less magic), \textbf{msgspec}
(extremely fast, schema-first, stricter), \textbf{dataclasses + manual
checks} (the honest baseline you now know how to beat).\\
\textbf{Installation.} \texttt{pip install "pydantic>=2.7"}.

\subsection{pydantic-settings}
\textbf{Description.} The \texttt{BaseSettings} layer: typed, environment-
layered configuration with prefixes, \texttt{.env} files, and explicit
precedence --- 12-factor factor III as a library~\cite{twelvefactor}.\\
\textbf{Why included.} The Orders API's \texttt{Settings} is its composition
root input; the test suite's determinism depends on injecting it explicitly.\\
\textbf{Alternatives.} \textbf{dynaconf} (multi-format, heavier),
\textbf{environs} (thin env parsing), \textbf{os.environ plus discipline}
(viable until the first time you need types or layering --- usually week
two).\\
\textbf{Installation.} \texttt{pip install "pydantic-settings>=2.2"}.

\subsection{httpx}
\textbf{Description.} The HTTP client whose \texttt{ASGITransport} powers
\texttt{TestClient}~\cite{httpxdocs}: requests are dispatched straight into
the ASGI callable, no socket involved.\\
\textbf{Why included.} Offline testing is this chapter's non-negotiable;
httpx is also the client library of Chapter 2, so one mental model serves
both directions of the wire.\\
\textbf{Alternatives.} \textbf{requests} (synchronous, no ASGI transport ---
fine for scripting, wrong here), \textbf{aiohttp} client (async, different
API).\\
\textbf{Installation.} \texttt{pip install "httpx>=0.27"}.

\subsection{pytest}
\textbf{Description.} The test runner: fixtures, parametrization, and the
\texttt{tmp\_path} fixture the SQLite test uses for hermetic file state.\\
\textbf{Why included.} The chapter's entire verification story --- 17 tests,
offline, one second --- is a pytest suite; Chapters 19 builds the full
quality-engineering practice on it.\\
\textbf{Alternatives.} \textbf{unittest} (stdlib, zero install, more
ceremony), \textbf{nose2} (largely historical).\\
\textbf{Installation.} \texttt{pip install "pytest>=8"}.

%% ----------------------------------------------------------------------------
\section{Chapter Summary \& Key Takeaways}
\label{sec:ch6-summary}

\begin{itemize}[itemsep=0.4em]
  \item \textbf{The contract is the product.} Design-first and code-first
        differ in who authors the OpenAPI document, not in whether it exists;
        pin, review, and CI-assert the generated schema and you get code-first
        speed with design-first governance.
  \item \textbf{ASGI is three callables.} \texttt{scope}, \texttt{receive},
        \texttt{send} --- uvicorn translates bytes to events, Starlette routes
        them, FastAPI compiles your annotations around them, and one event
        loop per worker runs the whole show.
  \item \textbf{\texttt{async} describes your I/O stack, not your wishes.}
        Any blocking link makes the endpoint \texttt{def}; the 40-thread
        pool is capacity you must size and instrument, as the threadpool
        exhaustion incident demonstrates.
  \item \textbf{DI is a per-request DAG with seams.} Sub-dependencies resolve
        depth-first, cache once per request, clean up after the response via
        \texttt{yield}, and swap wholesale via
        \texttt{dependency\_overrides}.
  \item \textbf{Pydantic v2 is the validation boundary and
        \texttt{response\_model} is the security boundary.} Field validators
        normalize, model validators enforce invariants, aliases decouple
        names, and the response model guarantees internal fields never reach
        the wire.
  \item \textbf{Operations are designed, not bolted on.} Settings through the
        environment, JSON logs with correlation IDs, three probes with three
        distinct meanings, and a lifespan that flips readiness before it
        closes anything.
  \item \textbf{Test the contract in-process.} \texttt{TestClient} drives the
        real app --- middleware, DI, validation, lifespan --- offline; assert
        status codes, error shapes, header echoes, and the OpenAPI document
        itself.
\end{itemize}

%% ----------------------------------------------------------------------------
\section{Quick-Reference Card}
\label{sec:ch6-quickref}

\subsection*{Endpoint Decision Table}
\begin{table}[h]
  \centering
  \small
  \begin{tabularx}{\textwidth}{l X}
    \toprule
    \textbf{Situation} & \textbf{Reach for} \\
    \midrule
    Async I/O end-to-end (async driver, \texttt{httpx}) & \texttt{async def} endpoint, awaited on the loop. \\
    Any blocking call (sync driver, \texttt{requests}) & \texttt{def} endpoint --- the threadpool absorbs it. \\
    CPU-bound work & Neither: subprocess or task queue; the loop and the pool are both wrong. \\
    Best-effort side effect after response & \texttt{BackgroundTasks}. \\
    Work that must survive a deploy & Durable task queue / outbox (Chapter 11). \\
    Constant-memory large payload & \texttt{StreamingResponse}; sync generator $=$ threadpool, async $=$ loop. \\
    \bottomrule
  \end{tabularx}
  \caption{Endpoint and execution-model decisions at a glance.}
\end{table}

\subsection*{Dependency-Injection Cheat-Sheet}
\begin{table}[h]
  \centering
  \small
  \begin{tabularx}{\textwidth}{l X}
    \toprule
    \textbf{Mechanism} & \textbf{Rule to remember} \\
    \midrule
    \texttt{Annotated[T, Depends(p)]} & Declare once as a type alias; reuse everywhere. \\
    Sub-dependencies & Providers can \texttt{Depends} too; the graph is flattened and resolved depth-first. \\
    Per-request cache & Each provider runs at most once per request; \texttt{use\_cache=False} opts out. \\
    \texttt{yield} cleanup & Post-\texttt{yield} runs after the response, LIFO, even on exceptions. \\
    Process-lifetime state & Build in \texttt{lifespan}, store on \texttt{app.state}, look up via providers. \\
    Test seam & \texttt{app.dependency\_overrides[original] = fake}; clear after the test. \\
    \bottomrule
  \end{tabularx}
  \caption{DI rules that prevent the common incidents.}
\end{table}

\subsection*{Lifecycle Ordering}
\begin{enumerate}
  \item Uvicorn parses bytes $\to$ builds the HTTP \texttt{scope}.
  \item Middleware, outermost first (last \texttt{add\_middleware} is outermost).
  \item Starlette router matches path and method.
  \item DI graph resolves (per-request cache populated).
  \item Pydantic validates inputs --- 422 short-circuits here.
  \item Handler runs: \texttt{async def} on the loop, \texttt{def} on the pool.
  \item Return value filtered through \texttt{response\_model}, serialized by alias.
  \item Middleware unwinds (headers already committed at
        \texttt{http.response.start}).
  \item \texttt{yield}-dependency cleanup runs (LIFO).
  \item Background tasks run; on SIGTERM the lifespan's post-\texttt{yield}
        code runs last.
\end{enumerate}

\section{Standardized Evidence Addendum}
\subsection{Benchmark Snapshot Template}
\begin{table}[h]
  \centering
  \small
  \begin{tabularx}{\textwidth}{l c c c c}
    \toprule
    \textbf{Config} & \textbf{p50 latency} & \textbf{p99 latency} & \textbf{Error rate} & \textbf{Throughput} \\
    \midrule
    Baseline & -- & -- & -- & 1.00x \\
    Candidate A & -- & -- & -- & -- \\
    Candidate B & -- & -- & -- & -- \\
    \bottomrule
  \end{tabularx}
  \caption{Minimum benchmark table to keep per chapter.}
\end{table}

\subsection{Request Trace Template}
\begin{itemize}
  \item \textbf{Request:} one representative API call for this chapter's topic.
  \item \textbf{Before:} behavior (headers, payload, latency, failure mode) with the naive approach.
  \item \textbf{After:} behavior with the technique in this chapter.
  \item \textbf{Result:} what changed in correctness, resilience, latency, and operability.
\end{itemize}

\subsection{Cost and Latency Budget Snapshot}
\begin{table}[h]
  \centering
  \small
  \begin{tabularx}{\textwidth}{l c c}
    \toprule
    \textbf{Stage} & \textbf{Latency budget} & \textbf{Cost driver} \\
    \midrule
    TLS / connection setup & -- & handshake round trips, keep-alive hit rate \\
    Request validation & -- & schema complexity, CPU per request \\
    Business logic and I/O & -- & downstream fan-out, query cost \\
    Serialization and egress & -- & payload size, compression ratio \\
    \bottomrule
  \end{tabularx}
  \caption{Operational budget template for this chapter's technique.}
\end{table}

\section{Configuration Preset Template}
\begin{minted}{text}
# api_policy.yaml
policy_version: "v1.0.0"
component: "<chapter_topic>"
timeout_ms: 0
retry_budget: 0
fallback_strategy: "fail_fast"
rollback_trigger: "error_budget_burn_gt_threshold"
\end{minted}

\section{CI Quality Gate Specification}
\begin{itemize}
  \item contract tests green against the pinned OpenAPI/AsyncAPI/Protobuf spec,
  \item no breaking schema changes without a version bump,
  \item error responses conform to RFC 9457 problem-details shape,
  \item benchmark deltas within approved regression bounds (latency and throughput),
  \item policy version and rollback plan present in release metadata.
\end{itemize}

\section{Incident Runbook Template}
\begin{enumerate}
  \item Detect regression from SLO burn-rate alerts or consumer reports.
  \item Scope impacted endpoints, versions, and consumer segments.
  \item Contain impact (rollback, feature flag, rate-limit tightening, or traffic shift).
  \item Diagnose with traces, structured logs, and contract diffs.
  \item Recover with validated fix and verified rollout procedure.
  \item Prevent recurrence via new regression tests and CI gates.
\end{enumerate}

\section{Cross-Chapter Decision Card}
\begin{table}[h]
  \centering
  \small
  \begin{tabularx}{\textwidth}{l X X}
    \toprule
    \textbf{Dimension} & \textbf{Choose this approach when} & \textbf{Prefer an alternative when} \\
    \midrule
    Use case & -- & -- \\
    Latency budget & -- & -- \\
    Operational cost & -- & -- \\
    Team maturity & -- & -- \\
    \bottomrule
  \end{tabularx}
  \caption{Decision card to complete per chapter.}
\end{table}

\section{ROI Model and Build-vs-Buy}
\begin{itemize}
  \item \textbf{Build when:} the capability is differentiating, compliance forbids vendors, or scale makes vendor pricing irrational.
  \item \textbf{Buy when:} the capability is commodity (auth, gateways, observability), time-to-market dominates, or staffing is thin.
  \item \textbf{Hidden costs to model:} on-call load, upgrade cadence, expertise ramp, data egress fees.
\end{itemize}

\section{Disaster Recovery Notes (RPO/RTO)}
\begin{itemize}
  \item Define recovery point objective (RPO): maximum acceptable data loss window.
  \item Define recovery time objective (RTO): maximum acceptable restoration time.
  \item Test restores regularly; an untested backup is not a backup.
\end{itemize}

